% Problems, Chapter 16: Angular Momentum
% Arfken, Weber & Harris, Mathematical Methods for Physicists, 7th ed.
% Requires: amsmath, amssymb, tikz

\section*{Problems}

Throughout this chapter use $\hbar=1$, as in the textbook. Scalar
spherical harmonics are normalized with the Condon--Shortley phase;
their indices satisfy $L\in\{0,1,\ldots\}$ and $-L\le M\le L$.
Terms beyond a ladder endpoint mean zero, rather than an undefined
harmonic. Abstract spin multiplets may also have half-integral $L$;
in that case $L\pm M$ and $2L$ are nonnegative integers.

\Needspace{12\baselineskip}
\subsection*{16.1 Angular Momentum Operators}

\Needspace{5\baselineskip}
\par\bigskip
\noindent\textbf{16.1.1}\quad The quantum mechanical angular momentum
operators $L_{x}\pm iL_{y}$ in 3-D physical space are given by
\begin{align*}
L_{x}+iL_{y}&=e^{i\varphi}
\left(\frac{\partial}{\partial\theta}
+i\cot\theta\,\frac{\partial}{\partial\varphi}\right),\\
L_{x}-iL_{y}&=-e^{-i\varphi}
\left(\frac{\partial}{\partial\theta}
-i\cot\theta\,\frac{\partial}{\partial\varphi}\right).
\end{align*}
Show that
\begin{enumerate}
\item[(a)] $(L_{x}+iL_{y})Y_{L}^{M}(\theta,\varphi)
=\sqrt{(L-M)(L+M+1)}\,Y_{L}^{M+1}(\theta,\varphi)$,
\item[(b)] $(L_{x}-iL_{y})Y_{L}^{M}(\theta,\varphi)
=\sqrt{(L+M)(L-M+1)}\,Y_{L}^{M-1}(\theta,\varphi)$.
\end{enumerate}
\solution
Use $\hbar=1$ and the Condon--Shortley normalization
\[
Y_L^M=N_{LM}P_L^M(\cos\theta)e^{iM\varphi},\qquad
N_{LM}=\sqrt{\frac{2L+1}{4\pi}\frac{(L-M)!}{(L+M)!}}.
\]
Put $P=P_L^M(\cos\theta)$,
$A=P_L^{M+1}(\cos\theta)$ and
$B=(L+M)(L-M+1)P_L^{M-1}(\cos\theta)$.
In the variable $x=\cos\theta$ the associated-Legendre recurrences
(15.87) and (15.91) read
\[
A+\frac{2Mx}{\sqrt{1-x^2}}\,P+B=0,\qquad
\sqrt{1-x^2}\,\frac{dP}{dx}=\frac{B-A}{2}.
\]
Since $\sqrt{1-x^2}=\sin\theta$, $x/\sqrt{1-x^2}=\cot\theta$ and
$d/d\theta=-\sin\theta\,d/dx$, they become
\[
-M\cot\theta\,P=\frac{A+B}{2},\qquad
\frac{dP}{d\theta}=\frac{A-B}{2}.
\]
Adding them, or subtracting the second from the first, gives
\[
(\partial_\theta-M\cot\theta)P=A,\qquad
(-\partial_\theta-M\cot\theta)P=B.
\]
The first identity also follows directly from the form
$P_L^M=c\,\sin^M\theta\,f(x)$ of Exercise 15.4.3, valid for
$-L\le M\le L$, where $c=(-1)^M/(2^LL!)$ and
$f=d^{L+M}(x^2-1)^L/dx^{L+M}$:
\[
\partial_\theta P
=M\cot\theta\,P-c\sin^{M+1}\theta\,f'(x)
=M\cot\theta\,P+P_L^{M+1},
\]
because the same form of $P_L^{M+1}$ has the constant
$(-1)^{M+1}/(2^LL!)=-c$ and the function $f'$.

Since $\partial_\varphi e^{iM\varphi}=iM e^{iM\varphi}$, the
azimuthal term of $L_+$ acts on $Y_L^M$ as
$i\cot\theta\cdot iM=-M\cot\theta$, so
\[
L_+Y_L^M=N_{LM}e^{i(M+1)\varphi}(\partial_\theta-M\cot\theta)P
=N_{LM}P_L^{M+1}e^{i(M+1)\varphi}
=\frac{N_{LM}}{N_{L,M+1}}Y_L^{M+1}.
\]
For $M<L$ the normalization ratio is
\[
\frac{N_{LM}^2}{N_{L,M+1}^2}
=\frac{(L-M)!(L+M+1)!}{(L+M)!(L-M-1)!}
=(L-M)(L+M+1),
\]
proving (a). Similarly, the azimuthal term of $L_-$ acts as
$-i\cot\theta\cdot iM=M\cot\theta$, so
\[
\begin{aligned}
L_-Y_L^M&=N_{LM}e^{i(M-1)\varphi}(-\partial_\theta-M\cot\theta)P
=N_{LM}B\,e^{i(M-1)\varphi}\\
&=(L+M)(L-M+1)
\frac{N_{LM}}{N_{L,M-1}}Y_L^{M-1},
\end{aligned}
\]
and, for $M>-L$,
\[
\frac{N_{LM}^2}{N_{L,M-1}^2}
=\frac{(L-M)!(L+M-1)!}{(L+M)!(L-M+1)!}
=\frac1{(L+M)(L-M+1)}.
\]
The coefficient of $Y_L^{M-1}$ is therefore
\[
\frac{(L+M)(L-M+1)}{\sqrt{(L+M)(L-M+1)}}=\sqrt{(L+M)(L-M+1)}.
\]
Taking positive square roots establishes
\[
\boxed{L_\pm Y_L^M=
\sqrt{(L\mp M)(L\pm M+1)}\,Y_L^{M\pm1}}.
\]
These formulas hold for integer $-L\le M\le L$ with the endpoint
understood as annihilation: $L_+Y_L^L=L_-Y_L^{-L}=0$.
Indeed $Y_L^{\pm L}$ is proportional to
$\sin^L\theta\,e^{\pm iL\varphi}$, and direct substitution gives zero:
\[
L_\pm(\sin^L\theta\,e^{\pm iL\varphi})
=\pm e^{\pm i(L+1)\varphi}
(L\sin^{L-1}\theta\cos\theta-L\cot\theta\,\sin^L\theta)=0.
\]
No harmonic outside the permitted $M$ range is needed.

\Needspace{5\baselineskip}
\par\bigskip
\noindent\textbf{16.1.2}\quad With $L_{\pm}$ given by
\[
L_{\pm}=L_{x}\pm iL_{y}=\pm e^{\pm i\varphi}
\left(\frac{\partial}{\partial\theta}
\pm i\cot\theta\,\frac{\partial}{\partial\varphi}\right),
\]
show that
\begin{enumerate}
\item[(a)] $Y_{l}^{m}=\sqrt{\dfrac{(l+m)!}{(2l)!\,(l-m)!}}\,
(L_{-})^{l-m}Y_{l}^{l}$,
\item[(b)] $Y_{l}^{m}=\sqrt{\dfrac{(l-m)!}{(2l)!\,(l+m)!}}\,
(L_{+})^{l+m}Y_{l}^{-l}$.
\end{enumerate}
\solution
Each application of $L_-$ changes $m$ to $m-1$ with the coefficient
proved in 16.1.1. Starting at $m=l$, step $j$ ($j=0,1,\ldots$) acts on
magnetic index $m=l-j$, so its coefficient is
$\sqrt{(l+m)(l-m+1)}=\sqrt{(2l-j)(j+1)}$. After $k$ steps we therefore
have
\[
(L_-)^kY_l^l
=\left[\prod_{j=0}^{k-1}(2l-j)(j+1)\right]^{1/2}Y_l^{l-k}.
\]
The two products are
$\prod_{j=0}^{k-1}(2l-j)=2l(2l-1)\cdots(2l-k+1)=(2l)!/(2l-k)!$ and
$\prod_{j=0}^{k-1}(j+1)=k!$. Set $k=l-m$, so that $2l-k=l+m$:
\[
(L_-)^{l-m}Y_l^l
=\sqrt{\frac{(2l)!(l-m)!}{(l+m)!}}\,Y_l^m.
\]
Dividing by this positive coefficient proves
\[
\boxed{Y_l^m=\sqrt{\frac{(l+m)!}{(2l)!(l-m)!}}\,
(L_-)^{l-m}Y_l^l}.
\]

For raising from $m=-l$, the successive squared coefficients are
again $(2l-j)(j+1)$, because at step $j$ the magnetic index is
$m=-l+j$ and $(l-m)(l+m+1)=(2l-j)(j+1)$.
Thus
\[
(L_+)^kY_l^{-l}
=\sqrt{\frac{(2l)!\,k!}{(2l-k)!}}Y_l^{-l+k}.
\]
Putting $k=l+m$, so that $2l-k=l-m$, gives
\[
(L_+)^{l+m}Y_l^{-l}
=\sqrt{\frac{(2l)!(l+m)!}{(l-m)!}}\,Y_l^m,
\]
and solving for the final harmonic gives
\[
\boxed{Y_l^m=\sqrt{\frac{(l-m)!}{(2l)!(l+m)!}}\,
(L_+)^{l+m}Y_l^{-l}}.
\]
For $k=0$ the product is empty and equals one, so both extreme
states are included.

\Needspace{5\baselineskip}
\par\bigskip
\noindent\textbf{16.1.3}\quad Using the known forms of $L_{+}$ and
$L_{-}$ (Exercise~16.1.2), show that
\[
\int\big[Y_{L}^{M}\big]^{*}\,L_{-}\big(L_{+}Y_{L}^{M}\big)\,d\Omega
=\int\big(L_{+}Y_{L}^{M}\big)^{*}\big(L_{+}Y_{L}^{M}\big)\,d\Omega.
\]
Here $d\Omega$ is the element of solid angle
($\sin\theta\,d\theta\,d\varphi$), and the integration is over the
entire angular space.
\solution
We prove the needed adjoint identity by integration by parts rather
than assuming it. Let $u,v$ be smooth single-valued functions on the
sphere, with $\varphi$-periodic coordinate expressions. Since
\[
L_-v=-e^{-i\varphi}v_\theta
+i e^{-i\varphi}\cot\theta\,v_\varphi,
\]
we have, with $d\Omega=\sin\theta\,d\theta\,d\varphi$ and
$\langle u,w\rangle=\int u^*w\,d\Omega$,
\[
\langle u,L_-v\rangle
=\int_0^{2\pi}\int_0^\pi
\left[-u^*e^{-i\varphi}\sin\theta\,v_\theta
+i u^*e^{-i\varphi}\cos\theta\,v_\varphi\right]d\theta\,d\varphi.
\]
Integrate the first term by parts in $\theta$ and the second in
$\varphi$:
\[
\begin{aligned}
-\int_0^\pi u^*e^{-i\varphi}\sin\theta\,v_\theta\,d\theta
&=\int_0^\pi v\,e^{-i\varphi}
(\sin\theta\,u_\theta^*+\cos\theta\,u^*)\,d\theta,\\
\int_0^{2\pi}iu^*e^{-i\varphi}\cos\theta\,v_\varphi\,d\varphi
&=-\int_0^{2\pi}v\,e^{-i\varphi}\cos\theta\,
(iu_\varphi^*+u^*)\,d\varphi,
\end{aligned}
\]
where $\partial_\varphi(iu^*e^{-i\varphi})
=e^{-i\varphi}(iu_\varphi^*+u^*)$.
The $\theta$ boundary term $[-u^*e^{-i\varphi}\sin\theta\,v]_0^\pi$
vanishes because $\sin\theta=0$ at the poles; the $\varphi$ boundary
term vanishes by periodicity. The remaining integrand is
\[
v\left[
e^{-i\varphi}\sin\theta\,u_\theta^*
+e^{-i\varphi}\cos\theta\,u^*
-e^{-i\varphi}\cos\theta\,u^*
-i e^{-i\varphi}\cos\theta\,u_\varphi^*
\right].
\]
The middle two terms cancel. Dividing the rest by $\sin\theta$ to
restore $d\Omega$, and using
$(L_+u)^*=e^{-i\varphi}(u_\theta^*-i\cot\theta\,u_\varphi^*)$,
we are left with
\[
\langle u,L_-v\rangle
=\int v\,e^{-i\varphi}(u_\theta^*
-i\cot\theta\,u_\varphi^*)\,d\Omega
=\langle L_+u,v\rangle.
\]
Take $u=Y_L^M$ and $v=L_+Y_L^M$, both regular spherical harmonics
or zero. This is exactly the requested equality. Moreover, since
$L_+Y_L^M=\sqrt{(L-M)(L+M+1)}\,Y_L^{M+1}$ by 16.1.1 and
$\int|Y_L^{M+1}|^2d\Omega=1$, both sides equal
$\boxed{(L-M)(L+M+1)}$, including zero at the upper endpoint.

\Needspace{5\baselineskip}
\par\bigskip
\noindent\textbf{16.1.4}
\begin{enumerate}
\item[(a)] Show that
$\mathbf{J}^{2}=\dfrac{1}{2}\big[J_{+}J_{-}+J_{-}J_{+}\big]+J_{z}^{2}$.
\item[(b)] Use the result from part (a) and the explicit formulas for
$L_{+}$ and $L_{-}$ from Exercise~16.1.2 to verify that all the
spherical harmonics with $l=2$ are eigenfunctions of $\mathbf{L}^{2}$
with eigenvalue $l(l+1)=6$.
\end{enumerate}
\solution
\textbf{(a)} Preserve operator order while expanding:
\[
\begin{aligned}
J_+J_-&=J_x^2+J_y^2+iJ_yJ_x-iJ_xJ_y,\\
J_-J_+&=J_x^2+J_y^2+iJ_xJ_y-iJ_yJ_x.
\end{aligned}
\]
Their sum is $2J_x^2+2J_y^2$ because the mixed terms cancel.
Consequently
\[
\boxed{\mathbf J^2=\tfrac12(J_+J_-+J_-J_+)+J_z^2}.
\]

\textbf{(b)} For an arbitrary smooth test function $f$, direct
product differentiation in the given ladder operators yields
\[
L_+L_-f=-f_{\theta\theta}-\cot\theta f_\theta
-\cot^2\theta f_{\varphi\varphi}-if_\varphi,
\]
\[
L_-L_+f=-f_{\theta\theta}-\cot\theta f_\theta
-\cot^2\theta f_{\varphi\varphi}+if_\varphi.
\]
For the first, write $c=\cot\theta$, so $c'=-\csc^2\theta$, and
$L_-f=-e^{-i\varphi}(f_\theta-icf_\varphi)$. Then
$L_+L_-f=e^{i\varphi}\partial_\theta(L_-f)
+e^{i\varphi}ic\,\partial_\varphi(L_-f)$, and since
$\partial_\varphi(-e^{-i\varphi})=ie^{-i\varphi}$ the two parts are
\[
\begin{aligned}
e^{i\varphi}\partial_\theta(L_-f)
&=-f_{\theta\theta}+ic'f_\varphi+ic f_{\theta\varphi},\\
e^{i\varphi}ic\,\partial_\varphi(L_-f)
&=ic\,[\,i(f_\theta-icf_\varphi)-(f_{\theta\varphi}-icf_{\varphi\varphi})\,]\\
&=-c f_\theta-ic f_{\theta\varphi}
+c^2if_\varphi-c^2f_{\varphi\varphi}.
\end{aligned}
\]
The mixed derivatives cancel and $i(c'+c^2)=-i$, because
$\csc^2\theta-\cot^2\theta=1$.
For the second, replacing $i$ by $-i$ in the coefficients turns
$L_+$ into $-L_-$ and $L_-$ into $-L_+$, hence turns $L_+L_-$ into
$(-L_-)(-L_+)=L_-L_+$. The computation above uses only $i^2=-1$ and
$\partial_\varphi e^{\pm i\varphi}=\pm ie^{\pm i\varphi}$, which remain
true with $-i$ in place of $i$, so the same replacement in its result
(with $f$ unchanged) gives $L_-L_+f$. This changes only the last
first-order term, from $-if_\varphi$ to $+if_\varphi$.
\erratum{The Instructor's Manual (printed p.~257) writes the difference
$(L_+L_--L_-L_+)/2$ when extracting the common second-order terms.
It must use the sum. The half-difference is $L_z=-i\partial_\varphi$,
not $\mathbf L^2-L_z^2$.} In the half-sum the terms $\mp if_\varphi$
cancel:
$\tfrac12(L_+L_-+L_-L_+)f=-f_{\theta\theta}-\cot\theta f_\theta
-\cot^2\theta f_{\varphi\varphi}$.
Since $L_z^2f=-f_{\varphi\varphi}$ and $\cot^2\theta+1=\csc^2\theta$,
(a) gives
\[
\mathbf L^2f=-f_{\theta\theta}-\cot\theta f_\theta
-\csc^2\theta f_{\varphi\varphi}.
\]
From here on write $s=\sin\theta$ and $c=\cos\theta$ (so $c$ no
longer denotes $\cot\theta$). The five $l=2$ harmonics are
nonzero constant multiples of
$s^2e^{\pm2i\varphi}$, $sc\,e^{\pm i\varphi}$, and $3c^2-1$.
The operator acts on $F(\theta)e^{im\varphi}$ as
$[-F''-\cot\theta F'+m^2F/s^2]e^{im\varphi}$. Explicitly,
\[
\begin{array}{c|c|c|c}
F&m^2&F',\,F''&-F''-(c/s)F'+m^2F/s^2\\ \hline
s^2&4&2sc,\ 2(c^2-s^2)&-2c^2+2s^2-2c^2+4=6s^2\\
sc&1&c^2-s^2,\ -4sc&4sc-c^3/s+cs+c/s=6sc\\
3c^2-1&0&-6sc,\ -6(c^2-s^2)&12c^2-6s^2=6(3c^2-1)
\end{array}
\]
The last column uses $s^2+c^2=1$: $4=4s^2+4c^2$ in the first row,
$c/s-c^3/s=c(1-c^2)/s=cs$ in the second, and $-6s^2=6c^2-6$ in the
third. Both signs of $m$ have the same $m^2$, so the first two rows each
verify two harmonics. Multiplying by their normalization constants
does not change the eigenvalue:
$\boxed{\mathbf L^2Y_2^m=6Y_2^m}$ for all five $m$.

\Needspace{5\baselineskip}
\par\bigskip
\noindent\textbf{16.1.5}\quad Derive the following relations for
\erratum{normalized eigenstates of the Hermitian angular-momentum
algebra in one irreducible multiplet, with relative phases chosen
so the ladder coefficients are positive}. Do not assume coordinate
formulas for $\psi_L^M$:
\begin{enumerate}
\item[(a)] $\psi_{L}^{M}(\theta,\varphi)
=\sqrt{\dfrac{(L+M)!}{(2L)!\,(L-M)!}}\,
(L_{-})^{L-M}\psi_{L}^{L}(\theta,\varphi)$,
\item[(b)] $\psi_{L}^{M}(\theta,\varphi)
=\sqrt{\dfrac{(L-M)!}{(2L)!\,(L+M)!}}\,
(L_{+})^{L+M}\psi_{L}^{-L}(\theta,\varphi)$.
\end{enumerate}
\solution
\erratum{The numerical coefficients require normalized states in a
single angular-momentum multiplet and a consistent ladder phase
convention. Eigenvalue equations alone leave arbitrary magnitudes,
phases, and possible degeneracy labels.}
Assume the Hermitian angular-momentum algebra, and choose states with
$\mathbf L^2\psi_L^M=L(L+1)\psi_L^M$ and
$L_z\psi_L^M=M\psi_L^M$. The commutators
$[L_z,L_\pm]=\pm L_\pm$ and $[\mathbf L^2,L_\pm]=0$ show that
$L_\pm\psi_L^M$ has the same $L$ and magnetic index $M\pm1$:
$L_zL_\pm\psi_L^M=L_\pm(L_z\pm1)\psi_L^M=(M\pm1)L_\pm\psi_L^M$.
Within a fixed irreducible multiplet it is therefore proportional
to $\psi_L^{M\pm1}$.

From $[L_x,L_y]=iL_z$,
$[L_+,L_-]=-i[L_x,L_y]+i[L_y,L_x]=2L_z$, while the identity of
16.1.4(a) gives $L_+L_-+L_-L_+=2(\mathbf L^2-L_z^2)$. Adding and
subtracting these two equations and halving gives
\[
L_-L_+=\mathbf L^2-L_z^2-L_z,\qquad
L_+L_-=\mathbf L^2-L_z^2+L_z.
\]
Since $L_\pm^\dagger=L_\mp$, normalization determines the squared
proportionality factors:
\[
\|L_+\psi_L^M\|^2=\langle\psi_L^M,L_-L_+\psi_L^M\rangle
=L(L+1)-M^2-M=(L-M)(L+M+1),
\]
\[
\|L_-\psi_L^M\|^2=\langle\psi_L^M,L_+L_-\psi_L^M\rangle
=L(L+1)-M^2+M=(L+M)(L-M+1).
\]
Choose the relative phases so these factors have positive square roots.
This works for raising and lowering at once: if
$L_-\psi_L^{M+1}=c\,\psi_L^M$ with $c>0$, then $L_+\psi_L^M$, being
proportional to $\psi_L^{M+1}$, has the coefficient
$\langle\psi_L^{M+1},L_+\psi_L^M\rangle
=\langle L_-\psi_L^{M+1},\psi_L^M\rangle=c$.
Thus fixing the phases of $\psi_L^{L-1},\ldots,\psi_L^{-L}$ by positive
lowering coefficients makes every raising coefficient positive too.
As in 16.1.2, lowering step $j$ acts on $M=L-j$, where
$(L+M)(L-M+1)=(2L-j)(j+1)$.
Iteration now gives, without any coordinate formula,
\[
(L_-)^k\psi_L^L
=\sqrt{\prod_{j=0}^{k-1}(2L-j)(j+1)}\,\psi_L^{L-k}
=\sqrt{\frac{(2L)!k!}{(2L-k)!}}\,\psi_L^{L-k}.
\]
Setting $k=L-M$ and inverting the coefficient yields (a).
Starting at $-L$ and raising gives the identical product; setting
$k=L+M$ yields (b):
\[
\boxed{\psi_L^M=\sqrt{\frac{(L+M)!}{(2L)!(L-M)!}}
(L_-)^{L-M}\psi_L^L},
\]
\[
\boxed{\psi_L^M=\sqrt{\frac{(L-M)!}{(2L)!(L+M)!}}
(L_+)^{L+M}\psi_L^{-L}}.
\]
The argument also applies to abstract half-integral multiplets:
$2L,L\pm M$ are still nonnegative integers. Ordinary scalar
spherical harmonics are the integral-$L$ realization.

\Needspace{5\baselineskip}
\par\bigskip
\noindent\textbf{16.1.6}\quad Derive the operator equations
\erratum{for integers $n\ge0$, on $0<\theta<\pi$, with the derivatives
in $\cos\theta$ taken at fixed $\varphi$}:
\begin{align*}
(L_{+})^{n}Y_{L}^{M}(\theta,\varphi)
&=(-1)^{n}e^{in\varphi}\sin^{n+M}\theta\,
\frac{d^{n}\sin^{-M}\theta\,Y_{L}^{M}(\theta,\varphi)}
{(d\cos\theta)^{n}},\\
(L_{-})^{n}Y_{L}^{M}(\theta,\varphi)
&=e^{-in\varphi}\sin^{n-M}\theta\,
\frac{d^{n}\sin^{M}\theta\,Y_{L}^{M}(\theta,\varphi)}
{(d\cos\theta)^{n}}.
\end{align*}
\par\emph{Hint.}~Try mathematical induction (Section~1.4).
\solution
All derivatives with respect to $x=\cos\theta$ below hold $\varphi$
fixed, so $\partial_\theta=-\sin\theta\,\partial_x$; the interior
range $0<\theta<\pi$ keeps $\cot\theta$ and the powers
$\sin^{\pm M}\theta$ finite. Since
$\partial_\theta\sin^{\mp m}\theta=\mp m\cot\theta\,\sin^{\mp m}\theta$,
a function $g_m=F(\theta)e^{im\varphi}$, for which
$\partial_\varphi g_m=img_m$, satisfies
\[
\begin{aligned}
L_+g_m
&=e^{i\varphi}(\partial_\theta-m\cot\theta)g_m
=e^{i\varphi}\sin^m\theta\,
\partial_\theta(\sin^{-m}\theta\,g_m)\\
&=-e^{i\varphi}\sin^{m+1}\theta\,
\partial_x(\sin^{-m}\theta\,g_m),\\
L_-g_m
&=-e^{-i\varphi}(\partial_\theta+m\cot\theta)g_m
=-e^{-i\varphi}\sin^{-m}\theta\,
\partial_\theta(\sin^m\theta\,g_m)\\
&=e^{-i\varphi}\sin^{1-m}\theta\,
\partial_x(\sin^m\theta\,g_m).
\end{aligned}
\]
These establish both claims for $n=1$; $n=0$ is the identity.

Assume the raising formula at $n$:
$g_{M+n}=(L_+)^nY_L^M=(-1)^ne^{in\varphi}\sin^{M+n}\theta
\,\partial_x^n(\sin^{-M}\theta\,Y_L^M)$.
It has magnetic index $M+n$, because $Y_L^M\propto e^{iM\varphi}$ and
$\partial_x$ at fixed $\varphi$ keeps that factor. Apply the first
identity with $m=M+n$. The two opposite powers of sine cancel inside
the derivative,
$\sin^{-(M+n)}\theta\,g_{M+n}
=(-1)^ne^{in\varphi}\partial_x^n(\sin^{-M}\theta\,Y_L^M)$,
and $e^{in\varphi}$ is constant in $x$, so
\[
L_+g_{M+n}
=(-1)^{n+1}e^{i(n+1)\varphi}\sin^{M+n+1}\theta\,
\partial_x^{n+1}(\sin^{-M}\theta\,Y_L^M).
\]
This is the $n+1$ formula.

For lowering, the induction hypothesis is
$g_{M-n}=(L_-)^nY_L^M=e^{-in\varphi}\sin^{n-M}\theta
\,\partial_x^n(\sin^M\theta\,Y_L^M)$.
Use the second identity with $m=M-n$; its factor
$\sin^m\theta$ cancels $\sin^{n-M}\theta$:
$\sin^{M-n}\theta\,g_{M-n}
=e^{-in\varphi}\partial_x^n(\sin^M\theta\,Y_L^M)$. Since
$\sin^{1-m}\theta=\sin^{n+1-M}\theta$, this leaves
\[
L_-g_{M-n}
=e^{-i(n+1)\varphi}\sin^{n+1-M}\theta\,
\partial_x^{n+1}(\sin^M\theta\,Y_L^M).
\]
Thus both stated equations hold for every integer $n\ge0$.
They are coordinate expressions on $0<\theta<\pi$; the resulting
harmonics extend regularly to the poles.

\Needspace{5\baselineskip}
\par\bigskip
\noindent\textbf{16.1.7}\quad Show, using $(L_{-})^{n}$, that
\[
Y_{L}^{-M}(\theta,\varphi)
=(-1)^{M}\big[Y_{L}^{M}(\theta,\varphi)\big]^{*}.
\]
\solution
First let $M\ge0$. Iterating the raising coefficient from zero to $M$
gives
\[
Y_L^M=c_M(L_+)^M Y_L^0,\qquad
c_M=\sqrt{\frac{(L-M)!}{(L+M)!}}.
\]
Indeed step $j$ acts on magnetic index $j$, with squared coefficient
$(L-j)(L+j+1)$ by 16.1.1, and the squared product is
\[
\prod_{j=0}^{M-1}(L-j)\prod_{j=0}^{M-1}(L+j+1)
=\frac{L!}{(L-M)!}\cdot\frac{(L+M)!}{L!}
=\frac{(L+M)!}{(L-M)!}.
\]
The differential forms show, for any differentiable $f$,
\[
(L_+f)^*
=e^{-i\varphi}(\partial_\theta-i\cot\theta\,\partial_\varphi)f^*
=-L_-f^*.
\]
Applying this $M$ times gives
$[(L_+)^Mf]^*=(-1)^M(L_-)^Mf^*$. Because $Y_L^0$ is real, therefore
\[
(Y_L^M)^*=(-1)^M c_M(L_-)^M Y_L^0.
\]
Lowering from zero to $-M$ has the same squared coefficient product
$(L+M)!/(L-M)!$, since step $j$ acts on $m=-j$ and
$(L+m)(L-m+1)=(L-j)(L+j+1)$. Hence $c_M(L_-)^M Y_L^0=Y_L^{-M}$, so
$(Y_L^M)^*=(-1)^MY_L^{-M}$. Multiplying by $(-1)^M$,
\[
\boxed{Y_L^{-M}=(-1)^M(Y_L^M)^*}.
\]
Replacing $M$ by $-M$ and conjugating proves the identical statement
for negative $M$, since $(-1)^{-M}=(-1)^M$ for integral $M$.
\erratum{The Instructor's Manual (printed p.~258) says to start with $Y_0^0$.
For a fixed degree $L$ the starting function must be $Y_L^0$.
Indeed both ladders annihilate $Y_0^0$, so they cannot generate
any nonconstant harmonic from it.}

\Needspace{5\baselineskip}
\par\bigskip
\noindent\textbf{16.1.8}\quad Verify by explicit calculation that
\begin{enumerate}
\item[(a)] $L_{+}Y_{1}^{0}(\theta,\varphi)
=-\sqrt{\dfrac{3}{4\pi}}\sin\theta\,e^{i\varphi}
=\sqrt{2}\,Y_{1}^{1}(\theta,\varphi)$,
\item[(b)] $L_{-}Y_{1}^{0}(\theta,\varphi)
=+\sqrt{\dfrac{3}{4\pi}}\sin\theta\,e^{-i\varphi}
=\sqrt{2}\,Y_{1}^{-1}(\theta,\varphi)$.
\end{enumerate}
The signs have the indicated values because the spherical harmonics
were defined to be consistent with the results obtained using the
ladder operators $L_{+}$ and $L_{-}$ (Condon--Shortley phase).
\solution
The relevant normalized harmonics are
\[
Y_1^0=\sqrt{\frac3{4\pi}}\cos\theta,\qquad
Y_1^1=-\sqrt{\frac3{8\pi}}\sin\theta\,e^{i\varphi},\qquad
Y_1^{-1}=\sqrt{\frac3{8\pi}}\sin\theta\,e^{-i\varphi}.
\]
Since $Y_1^0$ is independent of $\varphi$, the azimuthal derivatives
in $L_\pm$ vanish. Direct differentiation, with
$\partial_\theta\cos\theta=-\sin\theta$ and
$\sqrt{3/(4\pi)}=\sqrt2\sqrt{3/(8\pi)}$, gives
\[
L_+Y_1^0=e^{i\varphi}\sqrt{\frac3{4\pi}}(-\sin\theta)
=-\sqrt{\frac3{4\pi}}\sin\theta\,e^{i\varphi}
=\boxed{\sqrt2\,Y_1^1},
\]
\[
L_-Y_1^0=-e^{-i\varphi}\sqrt{\frac3{4\pi}}(-\sin\theta)
=\sqrt{\frac3{4\pi}}\sin\theta\,e^{-i\varphi}
=\boxed{\sqrt2\,Y_1^{-1}}.
\]
The minus sign in $Y_1^1$ is precisely the phase needed to make
both ladder coefficients positive.

\Needspace{12\baselineskip}
\subsection*{16.2 Angular Momentum Coupling}

\Needspace{5\baselineskip}
\par\bigskip
\noindent\textbf{16.2.1}\quad Derive recursion relations for
Clebsch--Gordan coefficients. Use them to calculate
$C(11J\,|\,m_{1}m_{2}M)$ for $J=0,1,2$.
\par\emph{Hint.}~Use the known matrix elements of
$J_{+}=J_{1+}+J_{2+}$, $J_{i+}$, and
$\mathbf{J}^{2}=(\mathbf{J}_{1}+\mathbf{J}_{2})^{2}$, etc.
\solution
Write $|m_1,m_2\rangle=|j_1m_1\rangle|j_2m_2\rangle$ and
\[
|JM\rangle=\sum_{m_1+m_2=M}C_M(m_1,m_2)|m_1,m_2\rangle.
\]
Define $a_i^\pm(m)=\sqrt{(j_i\mp m)(j_i\pm m+1)}$, so that
$J_{i\pm}|j_im\rangle=a_i^\pm(m)|j_i,m\pm1\rangle$.
Apply $J_+=J_{1+}+J_{2+}$ to this expansion and collect the
coefficient of the \emph{target} state $|m_1,m_2\rangle$
($m_1+m_2=M+1$); it comes from $|m_1-1,m_2\rangle$ through $J_{1+}$
and from $|m_1,m_2-1\rangle$ through $J_{2+}$, while the left side
is $\sqrt{(J-M)(J+M+1)}\,|J,M+1\rangle$:
\[
\begin{aligned}
\sqrt{(J-M)(J+M+1)}\,C_{M+1}(m_1,m_2)
={}&a_1^+(m_1-1)C_M(m_1-1,m_2)\\
&+a_2^+(m_2-1)C_M(m_1,m_2-1).
\end{aligned}
\]
Similarly, applying $J_-$ gives
\[
\begin{aligned}
\sqrt{(J+M)(J-M+1)}\,C_{M-1}(m_1,m_2)
={}&a_1^-(m_1+1)C_M(m_1+1,m_2)\\
&+a_2^-(m_2+1)C_M(m_1,m_2+1).
\end{aligned}
\]
Coefficients outside the allowed magnetic ranges are zero.
A recursion within a fixed $M$ follows from
\[
\mathbf J^2=\mathbf J_1^2+\mathbf J_2^2+2J_{1z}J_{2z}
+J_{1+}J_{2-}+J_{1-}J_{2+},
\]
since $J_{1+}J_{2-}+J_{1-}J_{2+}=2(J_{1x}J_{2x}+J_{1y}J_{2y})$, so the
last three terms are $2\mathbf J_1\cdot\mathbf J_2$. The first three
terms multiply $|m_1,m_2\rangle$ by
$j_1(j_1+1)+j_2(j_2+1)+2m_1m_2$; $J_{1+}J_{2-}$ reaches
$|m_1,m_2\rangle$ from $|m_1-1,m_2+1\rangle$ and $J_{1-}J_{2+}$ from
$|m_1+1,m_2-1\rangle$. Equating coefficients of $|m_1,m_2\rangle$ in
$\mathbf J^2|JM\rangle=J(J+1)|JM\rangle$ gives
\[
\begin{aligned}
&[J(J+1)-j_1(j_1+1)-j_2(j_2+1)-2m_1m_2]C_M(m_1,m_2)\\
&\quad=a_1^+(m_1-1)a_2^-(m_2+1)C_M(m_1-1,m_2+1)\\
&\qquad\quad+a_1^-(m_1+1)a_2^+(m_2-1)C_M(m_1+1,m_2-1).
\end{aligned}
\]
Together with normalization and phase, these determine the coefficients.

For $j_1=j_2=1$, choose positive lowering coefficients and a positive
coefficient of the largest $m_1$ in each highest-$M$ state, consistent
with the text's Condon--Shortley convention. For $j_i=1$,
$J_{i-}|1\rangle=\sqrt2|0\rangle$, $J_{i-}|0\rangle=\sqrt2|-1\rangle$
and $J_{i-}|-1\rangle=0$; the total ladder factors are
$J_-|J,M\rangle=\sqrt{(J+M)(J-M+1)}\,|J,M-1\rangle$.
Each lowering below applies the $J_-$ recursion to all coefficients
at once: the coefficient of $|m_1,m_2\rangle$ in $J_-|JM\rangle$ is
the right side of that recursion, and dividing by
$\sqrt{(J+M)(J-M+1)}$ gives $C_{M-1}(m_1,m_2)$.
Start with $|2,2\rangle=|1,1\rangle$. Since
$J_-|1,1\rangle=\sqrt2|0,1\rangle+\sqrt2|1,0\rangle$ and
$J_-|2,2\rangle=2|2,1\rangle$, the first lowering gives
\[
2|2,1\rangle
=\sqrt2(|1,0\rangle+|0,1\rangle).
\]
Next, $J_-(|1,0\rangle+|0,1\rangle)
=\sqrt2(|0,0\rangle+|1,-1\rangle)+\sqrt2(|-1,1\rangle+|0,0\rangle)$
and $J_-|2,1\rangle=\sqrt6|2,0\rangle$, so, after dividing by the
$\sqrt2$ in $|2,1\rangle$,
\[
\sqrt6|2,0\rangle
=|1,-1\rangle+2|0,0\rangle+|-1,1\rangle.
\]
Lowering this numerator gives
$\sqrt2|0,-1\rangle+2\sqrt2(|-1,0\rangle+|0,-1\rangle)
+\sqrt2|-1,0\rangle=3\sqrt2(|0,-1\rangle+|-1,0\rangle)$;
after dividing by its existing $\sqrt6$ and the next ladder factor
$\sqrt6$,
$|2,-1\rangle=3\sqrt2(|0,-1\rangle+|-1,0\rangle)/6
=(|0,-1\rangle+|-1,0\rangle)/\sqrt2$.
The final lowering gives
$2|2,-2\rangle=(\sqrt2+\sqrt2)|-1,-1\rangle/\sqrt2=2|-1,-1\rangle$.
Thus the complete quintet is
\[
\begin{array}{c|l}
M&|2,M\rangle\\ \hline
2&|1,1\rangle\\
1&(|1,0\rangle+|0,1\rangle)/\sqrt2\\
0&(|1,-1\rangle+2|0,0\rangle+|-1,1\rangle)/\sqrt6\\
-1&(|0,-1\rangle+|-1,0\rangle)/\sqrt2\\
-2&|-1,-1\rangle
\end{array}
\]

The normalized vector at $M=1$ orthogonal to $|2,1\rangle$, with a
positive coefficient of $|1,0\rangle$, is
$|1,1\rangle_{\rm coupled}=(|1,0\rangle-|0,1\rangle)/\sqrt2$.
Its two lowerings, with ladder factors
$\sqrt{(1+M)(2-M)}=\sqrt2$ for $M=1,0$, are
\[
\begin{aligned}
J_-|1,1\rangle_{\rm coupled}
&=\tfrac1{\sqrt2}\bigl[\sqrt2(|0,0\rangle+|1,-1\rangle)
-\sqrt2(|-1,1\rangle+|0,0\rangle)\bigr]\\
&=|1,-1\rangle-|-1,1\rangle=\sqrt2|1,0\rangle_{\rm coupled},
\end{aligned}
\]
\[
J_-|1,0\rangle_{\rm coupled}
=\tfrac1{\sqrt2}\bigl(\sqrt2|0,-1\rangle-\sqrt2|-1,0\rangle\bigr)
=|0,-1\rangle-|-1,0\rangle
=\sqrt2|1,-1\rangle_{\rm coupled}.
\]
Hence the triplet is
\[
\begin{array}{c|l}
M&|1,M\rangle_{\rm coupled}\\ \hline
1&(|1,0\rangle-|0,1\rangle)/\sqrt2\\
0&(|1,-1\rangle-|-1,1\rangle)/\sqrt2\\
-1&(|0,-1\rangle-|-1,0\rangle)/\sqrt2
\end{array}
\]
Finally write the singlet as
$u|1,-1\rangle+v|0,0\rangle+w|-1,1\rangle$.
Orthogonality to $|2,0\rangle$ and $|1,0\rangle_{\rm coupled}$ gives
$u+2v+w=0$ and $u-w=0$, so $w=u,v=-u$.
Normalization gives $3u^2=1$, and the chosen sign gives $u>0$:
\[
\boxed{|0,0\rangle_{\rm coupled}
=\frac{|1,-1\rangle-|0,0\rangle+|-1,1\rangle}{\sqrt3}}.
\]
The fixed-$M$ recursion checks all three $M=0$ states at once. For
$m_1=1$, $m_2=-1$ its bracket is $J(J+1)-2-2+2$, and only the term
with $a_1^+(0)a_2^-(0)=2$ survives on the right (the other needs
$m_1=2$), so
$C_0(0,0)=\tfrac12[J(J+1)-2]\,C_0(1,-1)$. This gives $2C_0(1,-1)$,
$0$ and $-C_0(1,-1)$ for $J=2,1,0$, as in the tables.

Every Clebsch--Gordan coefficient requested is the coefficient of its
product ket in these nine normalized states; all unlisted coefficients
are zero. Their dimensions $5+3+1=9$ exhaust the product basis.

\Needspace{5\baselineskip}
\par\bigskip
\noindent\textbf{16.2.2}\quad Defining $\big(Y_{l}\chi\big)_{J}^{M}$
by the formula
\[
\big(Y_{l}\chi\big)_{J}^{M}
=\sum C\big(l\tfrac{1}{2}J\,\big|\,m_{l}m_{s}M\big)
Y_{l}^{m_{l}}\chi_{m_{s}},
\]
where $\chi_{\pm1/2}$ are the spin up and down eigenfunctions of
$\sigma_{3}=\sigma_{z}$, show that $\big(Y_{l}\chi\big)_{J}^{M}$ is a
$J,M$ eigenfunction.
\solution
Let $\chi_+=\chi_{1/2}$ and $\chi_-=\chi_{-1/2}$, so that
$S_z\chi_\pm=\pm\chi_\pm/2$ with $\mathbf S=\boldsymbol\sigma/2$.
Every nonzero Clebsch--Gordan term has $m_l+m_s=M$; therefore
$J_z=L_z+S_z$ acting on the sum gives $M$ times that sum.

To check $\mathbf J^2$ explicitly, consider the fixed-$M$ product basis
\[
u=Y_l^{M-1/2}\chi_+,\qquad
v=Y_l^{M+1/2}\chi_-,
\]
omitting a vector if its orbital magnetic index is outside $[-l,l]$.
In this basis,
\[
\mathbf J^2=\mathbf L^2+\mathbf S^2+2L_zS_z+L_+S_-+L_-S_+
=\left[l(l+1)+\tfrac14\right]I+
\begin{pmatrix}M&q\\q&-M\end{pmatrix},
\]
where $K=l+\tfrac12$ and $q=\sqrt{K^2-M^2}$.
Indeed, by 16.1.1 and $S_-\chi_+=\chi_-$, $S_+\chi_-=\chi_+$,
\[
L_+S_-u=\sqrt{(l-M+\tfrac12)(l+M+\tfrac12)}\,v=\sqrt{(K-M)(K+M)}\,v=qv,
\]
and likewise $L_-S_+v=qu$, while $S_-\chi_-=S_+\chi_+=0$ make
$L_+S_-v=L_-S_+u=0$. With $\mathbf S^2=\tfrac34$ and
$2L_zS_z=2(M\mp\tfrac12)(\pm\tfrac12)$ on $u$ and $v$, respectively,
the diagonal terms are
$l(l+1)+\tfrac34\pm(M\mp\tfrac12)=l(l+1)+\tfrac14\pm M$.
The square of the last matrix is
$(M^2+q^2)I=K^2I$, so its eigenvalues are $\pm K$; its trace is zero,
so one is $+K$ and the other $-K$.

The corresponding normalized eigenvectors are
\[
\Psi_+=\sqrt{\frac{K+M}{2K}}\,u+
\sqrt{\frac{K-M}{2K}}\,v,
\qquad
\Psi_-=-\sqrt{\frac{K-M}{2K}}\,u+
\sqrt{\frac{K+M}{2K}}\,v.
\]
Indeed $q\sqrt{K-M}=(K-M)\sqrt{K+M}$ and
$q\sqrt{K+M}=(K+M)\sqrt{K-M}$. The two components of the matrix
times $\Psi_+$ are consequently
$[M+K-M]\sqrt{K+M}$ and $[K+M-M]\sqrt{K-M}$,
divided by $\sqrt{2K}$; these are the components of $K\Psi_+$.
For $\Psi_-$ the components are
$[-M+K+M]\sqrt{K-M}$ and
$[-K+M-M]\sqrt{K+M}$, divided by $\sqrt{2K}$,
which equal $-K\Psi_-$.
The complete eigenvalues are thus
\[
l(l+1)+\tfrac14+K=(l+\tfrac12)(l+\tfrac32),\qquad
l(l+1)+\tfrac14-K=(l-\tfrac12)(l+\tfrac12).
\]
These are $J(J+1)$ for $J=l+1/2$ and $J=l-1/2$.

The eigenvectors also carry the Condon--Shortley phases. Write
$u_M,v_M,\Psi_\pm(M)$ for the vectors above at magnetic index $M$.
Since $J_-=L_-+S_-$, 16.1.1 and $S_-\chi_+=\chi_-$, $S_-\chi_-=0$
give
\[
J_-u_M=\sqrt{(K+M-1)(K-M+1)}\,u_{M-1}+v_{M-1},\qquad
J_-v_M=q\,v_{M-1}.
\]
Hence, combining the $v_{M-1}$ terms with
$q\sqrt{K\mp M}=(K\mp M)\sqrt{K\pm M}$,
\[
\begin{aligned}
J_-\Psi_+(M)&=\sqrt{\tfrac{K+M}{2K}}\bigl[\sqrt{(K+M-1)(K-M+1)}\,u_{M-1}
+(K-M+1)\,v_{M-1}\bigr]\\
&=\sqrt{(K+M)(K-M+1)}\,\Psi_+(M-1),\\
J_-\Psi_-(M)&=\sqrt{\tfrac{K-M}{2K}}\bigl[-\sqrt{(K+M-1)(K-M+1)}\,u_{M-1}
+(K+M-1)\,v_{M-1}\bigr]\\
&=\sqrt{(K+M-1)(K-M)}\,\Psi_-(M-1).
\end{aligned}
\]
These are the positive ladder factors $\sqrt{(J+M)(J-M+1)}$ for $J=K$
and $J=K-1$. The top state of each multiplet has a positive
coefficient of $Y_l^l$: at $M=K$, $\Psi_+=u=Y_l^l\chi_+$, and at
$M=K-1$ the coefficient of $v=Y_l^l\chi_-$ in $\Psi_-$ is
$\sqrt{(2K-1)/2K}>0$. Hence $\Psi_\pm$ are the Condon--Shortley states,
and since $K\pm M=l\pm M+\tfrac12$ and $2K=2l+1$, their components are
the Clebsch--Gordan coefficients for $j_2=\tfrac12$:
\[
C(l\tfrac12J\,|\,M\mp\tfrac12,\pm\tfrac12,M)=
\begin{cases}
\sqrt{\dfrac{l\pm M+\frac12}{2l+1}},&J=l+\tfrac12,\\[2ex]
\mp\sqrt{\dfrac{l\mp M+\frac12}{2l+1}},&J=l-\tfrac12.
\end{cases}
\]
For $l=1$ these coefficients reproduce the states obtained by
explicit lowering in 16.2.3.
At the extreme $M=\pm(l+1/2)$ only the permitted basis vector
survives. The lower multiplet exists only for $l\ge1$.
Consequently
\[
\boxed{\mathbf J^2(Y_l\chi)_J^M=J(J+1)(Y_l\chi)_J^M,\qquad
J_z(Y_l\chi)_J^M=M(Y_l\chi)_J^M}.
\]

\Needspace{5\baselineskip}
\par\bigskip
\noindent\textbf{16.2.3}\quad Find the $(j,m)$ states of a $p$
electron ($l=1$), in which the orbital angular momentum of the
electron is coupled to its spin angular momentum ($s=1/2$) to form
states whose conventional labelings are ${}^{2}p_{1/2}$ and
${}^{2}p_{3/2}$. The notation is of the general form
${}^{2s+1}(\mathrm{symbol})_{j}$, where ``symbol'' is that indicating
the $l$ value (i.e., $s$, $p$, \ldots).
\solution
Write $p_m=Y_1^m$ and $\alpha=\chi_{1/2}$,
$\beta=\chi_{-1/2}$. The tensor products are normalized and mutually
orthogonal. Since $1\otimes\tfrac12=\tfrac32\oplus\tfrac12$,
there are four \( {}^{2}p_{3/2}\) states and two \( {}^{2}p_{1/2}\) states.

The highest state is $|\tfrac32,\tfrac32\rangle=p_1\alpha$.
Use $L_-p_1=\sqrt2p_0$, $L_-p_0=\sqrt2p_{-1}$,
$S_-\alpha=\beta$, $S_-\beta=0$:
\[
J_-(p_1\alpha)=\sqrt2p_0\alpha+p_1\beta.
\]
Divide by $\sqrt3$, the $J=3/2$ ladder coefficient
$\sqrt{(J+M)(J-M+1)}$ at $M=3/2$, to get
$|\tfrac32,\tfrac12\rangle$ below. Lower its numerator once more:
\[
J_-(\sqrt2p_0\alpha+p_1\beta)
=\sqrt2(\sqrt2p_{-1}\alpha+p_0\beta)+\sqrt2p_0\beta
=2p_{-1}\alpha+2\sqrt2p_0\beta.
\]
Dividing by the previous $\sqrt3$ and the next ladder coefficient
$\sqrt{2\cdot2}=2$ gives $|\tfrac32,-\tfrac12\rangle
=(p_{-1}\alpha+\sqrt2p_0\beta)/\sqrt3$.
The last lowering is
$J_-(p_{-1}\alpha+\sqrt2p_0\beta)/\sqrt3
=(p_{-1}\beta+2p_{-1}\beta)/\sqrt3=\sqrt3p_{-1}\beta$, and the ladder
coefficient at $M=-1/2$ is $\sqrt{1\cdot3}=\sqrt3$. Therefore
\[
\boxed{\begin{aligned}
|\tfrac32,\tfrac32\rangle&=p_1\alpha,\\
|\tfrac32,\tfrac12\rangle&=\sqrt{\tfrac23}p_0\alpha+
\sqrt{\tfrac13}p_1\beta,\\
|\tfrac32,-\tfrac12\rangle&=\sqrt{\tfrac13}p_{-1}\alpha+
\sqrt{\tfrac23}p_0\beta,\\
|\tfrac32,-\tfrac32\rangle&=p_{-1}\beta.
\end{aligned}}
\]
For the remaining $M=1/2$ state, orthogonality to
$|\tfrac32,\tfrac12\rangle$ requires $\sqrt2u+v=0$ in
$u p_0\alpha+v p_1\beta$. Normalization gives $|u|^2+|v|^2=3u^2=1$;
choosing the coefficient of $p_1\beta$ ($m_l=l$) positive, as in
16.2.1, fixes $u=-1/\sqrt3$, $v=\sqrt{2/3}$.
Its lowering coefficient for $J=1/2$ is $\sqrt{1\cdot1}=1$:
\[
\begin{aligned}
J_-\left(-\frac{p_0\alpha}{\sqrt3}
+\sqrt{\frac23}p_1\beta\right)
&=-\sqrt{\frac23}p_{-1}\alpha
-\frac{p_0\beta}{\sqrt3}+\frac2{\sqrt3}p_0\beta\\
&=-\sqrt{\frac23}p_{-1}\alpha+\frac{p_0\beta}{\sqrt3}.
\end{aligned}
\]
Thus
\[
\boxed{\begin{aligned}
|\tfrac12,\tfrac12\rangle&=-\sqrt{\tfrac13}p_0\alpha+
\sqrt{\tfrac23}p_1\beta,\\
|\tfrac12,-\tfrac12\rangle&=-\sqrt{\tfrac23}p_{-1}\alpha+
\sqrt{\tfrac13}p_0\beta.
\end{aligned}}
\]
\erratum{The Instructor's Manual's displayed doublet changes the relative phase
between its two $M$ states and does not obey
$J_-|\tfrac12,\tfrac12\rangle=|\tfrac12,-\tfrac12\rangle$.
The explicit lowering above fixes that phase. A common overall sign
change of both states would be harmless.}
The superscript $2$ is $2s+1$, and the letter $p$ records $l=1$.

\Needspace{5\baselineskip}
\par\bigskip
\noindent\textbf{16.2.4}\quad Repeat Exercise~16.2.3 for $l=1$,
$s=3/2$. Apply the conventional labels to the $j,m$ states.
\solution
Here $1\otimes\tfrac32=\tfrac52\oplus\tfrac32\oplus\tfrac12$.
Because $2s+1=4$ and $l=1$, the labels are
\( {}^4p_{5/2}\), \( {}^4p_{3/2}\), and \( {}^4p_{1/2}\),
with $6,4,2$ states, respectively.
Write $p_+=Y_1^1$, $p_0=Y_1^0$, $p_-=Y_1^{-1}$ and let
$\chi_\mu$ be the normalized spin-$3/2$ state with $S_z\chi_\mu=\mu\chi_\mu$.
The elementary lowering actions are
\[
L_-p_+=\sqrt2p_0,\quad L_-p_0=\sqrt2p_-,\quad L_-p_-=0,
\]
\[
S_-\chi_{3/2}=\sqrt3\chi_{1/2},\quad
S_-\chi_{1/2}=2\chi_{-1/2},\quad
S_-\chi_{-1/2}=\sqrt3\chi_{-3/2},\quad S_-\chi_{-3/2}=0.
\]
These specify every matrix element used below.

For a compact complete table, express a state of given $M$ by its
coefficient column in the ordered product basis
\[
B_M=(p_+\chi_{M-1},\ p_0\chi_M,\ p_-\chi_{M+1}),
\qquad |JM\rangle=B_Mc_M^{(J)}.
\]
A nonexistent spin state is omitted, and its coefficient is zero.
Let $\sigma(\mu)=\sqrt{(3/2+\mu)(5/2-\mu)}$ for allowed $\mu$;
set $\sigma=0$ outside the spin range. Thus
$\sigma(\tfrac32)=\sqrt3$, $\sigma(\tfrac12)=2$,
$\sigma(-\tfrac12)=\sqrt3$, $\sigma(-\tfrac32)=0$, the $S_-$
coefficients above. Applying $J_-=L_-+S_-$
and collecting the three target coefficients gives the explicit rule
\[
\sqrt{(J+M)(J-M+1)}\,c_{M-1}^{(J)}
=
\begin{pmatrix}
\sigma(M-1)&0&0\\
\sqrt2&\sigma(M)&0\\
0&\sqrt2&\sigma(M+1)
\end{pmatrix}c_M^{(J)}.
\tag{*}
\]
Its columns are the images of the three basis states:
$J_-(p_+\chi_{M-1})=\sigma(M-1)p_+\chi_{M-2}+\sqrt2p_0\chi_{M-1}$,
$J_-(p_0\chi_M)=\sigma(M)p_0\chi_{M-1}+\sqrt2p_-\chi_M$, and
$J_-(p_-\chi_{M+1})=\sigma(M+1)p_-\chi_M$.
Each row of the tables below lists the complete coefficient column
$c_M^{(J)}$ of $|JM\rangle$ in the basis $B_M$.

For $J=5/2$, start with $c_{5/2}=(1,0,0)^T$.
The first lowering gives $(\sigma(\tfrac32),\sqrt2,0)^T
=(\sqrt3,\sqrt2,0)^T$, whose norm $\sqrt5$ equals the ladder factor
$\sqrt{5\cdot1}$. The next, with $\sigma(\tfrac12)=2$,
$\sigma(\tfrac32)=\sqrt3$, gives
\[
\frac1{\sqrt5}(2\sqrt3,\ \sqrt2\sqrt3+\sqrt3\sqrt2,\ \sqrt2\sqrt2)^T
=\frac1{\sqrt5}(2\sqrt3,\ 2\sqrt6,\ 2)^T;
\]
dividing by the next factor $\sqrt{4\cdot2}=2\sqrt2$ gives the
$M=1/2$ row.
For the remaining steps the factors in (*) are
$\sqrt{3\cdot3}=3$, $\sqrt{2\cdot4}=2\sqrt2$, $\sqrt{1\cdot5}=\sqrt5$,
and the numerators are ($\sigma(-\tfrac32)=0$ in the last two)
\[
\begin{aligned}
M=\tfrac12:&\quad
\bigl(\sqrt3\sqrt{\tfrac3{10}},\
\sqrt2\sqrt{\tfrac3{10}}+2\sqrt{\tfrac35},\
\sqrt2\sqrt{\tfrac35}+\sqrt3\sqrt{\tfrac1{10}}\bigr)^T
=\bigl(\tfrac3{\sqrt{10}},\ 3\sqrt{\tfrac35},\
3\sqrt{\tfrac3{10}}\bigr)^T,\\
M=-\tfrac12:&\quad
\bigl(0,\ \sqrt2\sqrt{\tfrac1{10}}+\sqrt3\sqrt{\tfrac35},\
\sqrt2\sqrt{\tfrac35}+2\sqrt{\tfrac3{10}}\bigr)^T
=\bigl(0,\ \tfrac4{\sqrt5},\ 2\sqrt{\tfrac65}\bigr)^T,\\
M=-\tfrac32:&\quad
\bigl(0,\ 0,\ \sqrt2\sqrt{\tfrac25}+\sqrt3\sqrt{\tfrac35}\bigr)^T
=(0,\ 0,\ \sqrt5)^T,
\end{aligned}
\]
where $\sqrt2\sqrt{3/10}=\sqrt{3/5}$,
$\sqrt2\sqrt{3/5}=\sqrt{6/5}=2\sqrt{3/10}$,
$\sqrt3\sqrt{1/10}=\sqrt{3/10}$ and $\sqrt2\sqrt{1/10}=1/\sqrt5$.
Dividing by the factors $3$, $2\sqrt2$, $\sqrt5$ gives the last three
rows of the complete sextet:
\[
\boxed{\begin{array}{c|ccc}
M&c_+&c_0&c_-\\ \hline
5/2&1&0&0\\
3/2&\sqrt{3/5}&\sqrt{2/5}&0\\
1/2&\sqrt{3/10}&\sqrt{3/5}&\sqrt{1/10}\\
-1/2&\sqrt{1/10}&\sqrt{3/5}&\sqrt{3/10}\\
-3/2&0&\sqrt{2/5}&\sqrt{3/5}\\
-5/2&0&0&1
\end{array}}\qquad ({}^4p_{5/2}).
\]

At $M=3/2$, the highest state of the $J=3/2$ multiplet must be
orthogonal to $(\sqrt{3/5},\sqrt{2/5},0)^T$. A positive first
coefficient ($m_l=l$, as in 16.2.1) fixes it to
$(\sqrt{2/5},-\sqrt{3/5},0)^T$.
Its three lowerings, with factors $\sqrt{3\cdot1}=\sqrt3$,
$\sqrt{2\cdot2}=2$ and $\sqrt{1\cdot3}=\sqrt3$, have numerators
\[
\begin{aligned}
M=\tfrac32:&\quad
\bigl(2\sqrt{\tfrac25},\ \sqrt2\sqrt{\tfrac25}-\sqrt3\sqrt{\tfrac35},\
-\sqrt2\sqrt{\tfrac35}\bigr)^T
=\bigl(2\sqrt{\tfrac25},\ -\tfrac1{\sqrt5},\ -\sqrt{\tfrac65}\bigr)^T,\\
M=\tfrac12:&\quad
\bigl(\sqrt3\sqrt{\tfrac8{15}},\
\sqrt2\sqrt{\tfrac8{15}}-\tfrac2{\sqrt{15}},\
-\tfrac{\sqrt2}{\sqrt{15}}-\sqrt3\sqrt{\tfrac25}\bigr)^T
=\bigl(2\sqrt{\tfrac25},\ \tfrac2{\sqrt{15}},\
-2\sqrt{\tfrac8{15}}\bigr)^T,\\
M=-\tfrac12:&\quad
\bigl(0,\ \sqrt2\sqrt{\tfrac25}+\tfrac{\sqrt3}{\sqrt{15}},\
\tfrac{\sqrt2}{\sqrt{15}}-2\sqrt{\tfrac8{15}}\bigr)^T
=\bigl(0,\ \tfrac3{\sqrt5},\ -\sqrt{\tfrac65}\bigr)^T,
\end{aligned}
\]
where $\sqrt2\sqrt{2/5}=2/\sqrt5$, $\sqrt3\sqrt{3/5}=3/\sqrt5$,
$\sqrt2\sqrt{8/15}=4/\sqrt{15}$, $\sqrt3\sqrt{2/5}=3\sqrt{2/15}$ and
$2\sqrt{8/15}=4\sqrt{2/15}$; for example the last component is
$\sqrt{2/15}-4\sqrt{2/15}=-3\sqrt{2/15}=-\sqrt{6/5}$.
Dividing by the factors gives
\[
\boxed{\begin{array}{c|ccc}
M&c_+&c_0&c_-\\ \hline
3/2&\sqrt{2/5}&-\sqrt{3/5}&0\\
1/2&\sqrt{8/15}&-1/\sqrt{15}&-\sqrt{2/5}\\
-1/2&\sqrt{2/5}&1/\sqrt{15}&-\sqrt{8/15}\\
-3/2&0&\sqrt{3/5}&-\sqrt{2/5}
\end{array}}\qquad ({}^4p_{3/2}).
\]

For the remaining highest state at $M=1/2$, write
$c=(A,B,C)^T$. Orthogonality to the two existing $M=1/2$ rows
(multiplied by $\sqrt{10}$ and $\sqrt{15}$, respectively) gives
\[
\sqrt3 A+\sqrt6 B+C=0,\qquad
2\sqrt2 A-B-\sqrt6 C=0.
\]
Substituting $C=-\sqrt3 A-\sqrt6 B$ in the second equation yields
$2\sqrt2A-B+3\sqrt2A+6B=5\sqrt2 A+5B=0$, so $B=-\sqrt2 A$ and
$C=-\sqrt3A+2\sqrt3A=\sqrt3 A$.
Normalization gives $(1+2+3)A^2=6A^2=1$; choose $A=1/\sqrt6>0$.
Then $c_{1/2}=(1/\sqrt6,-1/\sqrt3,1/\sqrt2)^T$.
Applying (*) with $\sigma(-\tfrac12)=\sqrt3$, $\sigma(\tfrac12)=2$,
$\sigma(\tfrac32)=\sqrt3$ and its factor $\sqrt{1\cdot1}=1$ produces
\[
\left(\frac1{\sqrt2},\
\frac1{\sqrt3}-\frac2{\sqrt3},\
-\sqrt{\frac23}+\sqrt{\frac32}\right)^T
=\left(\frac1{\sqrt2},-\frac1{\sqrt3},\frac1{\sqrt6}\right)^T.
\]
Thus the doublet is
\[
\boxed{\begin{array}{c|ccc}
M&c_+&c_0&c_-\\ \hline
1/2&1/\sqrt6&-1/\sqrt3&1/\sqrt2\\
-1/2&1/\sqrt2&-1/\sqrt3&1/\sqrt6
\end{array}}\qquad ({}^4p_{1/2}).
\]
Each row has norm one; rows of the same $M$ but different $J$ are
orthogonal, being eigenvectors of the Hermitian operator $\mathbf J^2$
with different eigenvalues. For instance, at $M=1/2$ the $J=5/2$ and
$J=3/2$ rows give
$\sqrt{3/10}\sqrt{8/15}-\sqrt{3/5}/\sqrt{15}-\sqrt{1/10}\sqrt{2/5}
=(2-1-1)/5=0$. The dimension
$6+4+2=12=3\cdot4$ confirms that every product-basis state is included.
A common sign change within a complete multiplet changes no physics;
the tables use the same ladder phase convention as 16.2.1.
\erratum{The Instructor's Manual's $J=3/2,M=-1/2$ row contains an invalid spin
label on its $p_+$ term. That label must be $-3/2$, since
$1+(-3/2)=-1/2$. The corrected table above obeys $m_l+m_s=M$
in every entry. Its $J=3/2$ multiplet differs from the Instructor's Manual by
one common, harmless overall sign. Its intermediate lowering
work on printed p.~261 also has two $p_0$ labels that must be $p_-$:
the term with spin $-1/2$ in $J=5/2,M=-3/2$, and the term with
spin $+1/2$ in $J=1/2,M=-1/2$. The sums $m_l+m_s$ fix both labels.}

\Needspace{5\baselineskip}
\par\bigskip
\noindent\textbf{16.2.5}\quad A deuterium atom consists of a proton,
a neutron, and an electron. Each of these particles has spin $1/2$.
The coupling of these three spins can produce $J$ values of $3/2$ and
$1/2$. We consider here only states with no orbital angular momentum.
\begin{enumerate}
\item[(a)] Show that these $J,M$ states consist of one quartet
($J=3/2$) and two linearly independent doublets ($J=1/2$).
\par\emph{Hint.}~Make a vector-model diagram.
\item[(b)] One way to analyze this problem is to couple the spins of
the proton and neutron to form a nuclear triplet or singlet, and then
to couple the resultant nuclear spin to the electron spin. Find the
states that are obtained in this way (designate the single-particle
states $p_{\alpha}$, $p_{\beta}$, $n_{\alpha}$, $n_{\beta}$,
$e_{\alpha}$, $e_{\beta}$).
\item[(c)] Another way to analyze this problem is to couple the spins
of the proton and electron to form an atomic triplet or singlet, and
then to couple that resultant to the neutron spin. Find the states
that result from this coupling scheme.
\item[(d)] Show that the coupling schemes of parts (b) and (c) span
the same Hilbert space.
\end{enumerate}
\par\emph{Note.}~The actual interaction energies among these angular
momenta cause the scheme of part (b) to be the better way of treating
this problem (the triplet nuclear state is substantially the more
stable), and the system actually looks like a spin-1 deuterium
nucleus plus an electron.
\solution
Use the fixed tensor order proton--neutron--electron.
The particles are distinct, so reordering their tensor factors while
writing a coupling scheme introduces no fermionic exchange sign.

\textbf{(a)} The product space has dimension $2^3=8$. Its numbers
of states at $M=3/2,1/2,-1/2,-3/2$ are respectively $1,3,3,1$,
the numbers $\binom30,\binom31,\binom32,\binom33$ of ways to choose
which of the three spins point down.
The unique highest state generates one $J=3/2$ quartet, removing
one state from each row. The remainder has two states at each of
$M=\pm1/2$, hence two independent $J=1/2$ doublets.
Equivalently, the coupling tree is
\[
(\tfrac12\otimes\tfrac12)\otimes\tfrac12
=(1\oplus0)\otimes\tfrac12
=(\tfrac32\oplus\tfrac12)\oplus\tfrac12.
\]

\textbf{(b)} Couple $p$ and $n$ first. Their triplet and singlet are
\[
t_+=p_\alpha n_\alpha,\qquad
t_0=\frac{p_\alpha n_\beta+p_\beta n_\alpha}{\sqrt2},\qquad
t_-=p_\beta n_\beta,\qquad
s=\frac{p_\alpha n_\beta-p_\beta n_\alpha}{\sqrt2}.
\]
With $J_-=J_{p-}+J_{n-}$ on the nuclear states, lowering each spin in
turn gives $J_-t_+=p_\beta n_\alpha+p_\alpha n_\beta=\sqrt2\,t_0$ and
$J_-t_0=2p_\beta n_\beta/\sqrt2=\sqrt2\,t_-$, while $J_\pm s=0$; so
the $t_m$ have the ladder coefficients of $p_m=Y_1^m$ in 16.2.3.
Applying the $1\otimes1/2$ formulas of 16.2.3 with $p_m\to t_m$,
$\alpha\to e_\alpha$, $\beta\to e_\beta$ gives the quartet
\[
\begin{aligned}
Q_{3/2}&=t_+e_\alpha=p_\alpha n_\alpha e_\alpha,\\
Q_{1/2}&=\sqrt{\tfrac23}t_0e_\alpha+\sqrt{\tfrac13}t_+e_\beta,\\
Q_{-1/2}&=\sqrt{\tfrac13}t_-e_\alpha+\sqrt{\tfrac23}t_0e_\beta,\\
Q_{-3/2}&=t_-e_\beta=p_\beta n_\beta e_\beta,
\end{aligned}
\]
and two doublets, denoted $T_M$ for nuclear triplet and $S_M$ for
nuclear singlet:
\[
\begin{aligned}
T_{1/2}&=\sqrt{\tfrac23}t_+e_\beta-\sqrt{\tfrac13}t_0e_\alpha,&
S_{1/2}&=s e_\alpha,\\
T_{-1/2}&=\sqrt{\tfrac13}t_0e_\beta-\sqrt{\tfrac23}t_-e_\alpha,&
S_{-1/2}&=s e_\beta.
\end{aligned}
\]
To display all product coefficients and compare the next scheme, define
\[
\begin{array}{lll}
A=p_\alpha n_\alpha e_\beta,&
B=p_\alpha n_\beta e_\alpha,&
C=p_\beta n_\alpha e_\alpha,\\
D=p_\alpha n_\beta e_\beta,&
E=p_\beta n_\alpha e_\beta,&
F=p_\beta n_\beta e_\alpha.
\end{array}
\]
Substitution of $t_0,s$ then gives
\[
\begin{array}{c|ccc}
M&Q_M&T_M&S_M\\ \hline
1/2&(A+B+C)/\sqrt3&(2A-B-C)/\sqrt6&(B-C)/\sqrt2\\
-1/2&(D+E+F)/\sqrt3&(D+E-2F)/\sqrt6&(D-E)/\sqrt2
\end{array}
\]
together with the two extreme quartet states already written.
For example $t_+e_\beta=A$ and $t_0e_\alpha=(B+C)/\sqrt2$, so
$T_{1/2}=\sqrt{2/3}\,A-\sqrt{1/3}\,(B+C)/\sqrt2=(2A-B-C)/\sqrt6$;
the other entries follow in the same way.
Lowering each spin in turn gives $J_-A=D+E$, $J_-B=D+F$ and
$J_-C=E+F$ (for example, in $A=p_\alpha n_\alpha e_\beta$ lowering
$p$ gives $E$, lowering $n$ gives $D$, and $e_\beta$ is annihilated).
Consequently $J_-(2A-B-C)=D+E-2F$ and
$J_-(B-C)=D-E$. Since the $J=1/2$ ladder coefficient at $M=1/2$ is
$\sqrt{1\cdot1}=1$, this explicitly checks the doublet ladder phases.

\textbf{(c)} Replace the first pair by $pe$:
\[
t'_+=p_\alpha e_\alpha,\quad
t'_0=(p_\alpha e_\beta+p_\beta e_\alpha)/\sqrt2,\quad
t'_-=p_\beta e_\beta,\quad
s'=(p_\alpha e_\beta-p_\beta e_\alpha)/\sqrt2.
\]
Couple these to $n$ with the same coefficients:
\[
\begin{aligned}
T'_{1/2}&=\sqrt{\tfrac23}t'_+n_\beta-\sqrt{\tfrac13}t'_0n_\alpha,&
S'_{1/2}&=s'n_\alpha,\\
T'_{-1/2}&=\sqrt{\tfrac13}t'_0n_\beta-\sqrt{\tfrac23}t'_-n_\alpha,&
S'_{-1/2}&=s'n_\beta.
\end{aligned}
\]
Reordering the displayed factors back to $pne$ (for example
$t'_+n_\beta=p_\alpha n_\beta e_\alpha=B$ and
$t'_0n_\alpha=(A+C)/\sqrt2$, so
$T'_{1/2}=\sqrt{2/3}\,B-\sqrt{1/3}\,(A+C)/\sqrt2$) gives
\[
\begin{array}{c|cc}
M&T'_M&S'_M\\ \hline
1/2&(2B-A-C)/\sqrt6&(A-C)/\sqrt2\\
-1/2&(D-2E+F)/\sqrt6&(D-F)/\sqrt2
\end{array}
\]
and the same four $Q_M$: the quartet is totally symmetric under
permuting its spin labels.

\textbf{(d)} At each $M=\pm1/2$ the product space is
three-dimensional, with orthonormal basis $Q_M,T_M,S_M$. The states
$T'_M,S'_M$ have $J=1/2$, so they are orthogonal to the $J=3/2$
state $Q_M$ (eigenvectors of the Hermitian $\mathbf J^2$ with different
eigenvalues), and therefore
$T'_M=\langle T_M,T'_M\rangle T_M+\langle S_M,T'_M\rangle S_M$, and
similarly for $S'_M$. For $M=1/2$ the coefficient vectors in the
basis $(A,B,C)$ are $T_M=(2,-1,-1)/\sqrt6$, $S_M=(0,1,-1)/\sqrt2$,
$T'_M=(-1,2,-1)/\sqrt6$, $S'_M=(1,0,-1)/\sqrt2$, and their dot
products give
\[
\langle T_M,T'_M\rangle=(-2-2+1)/6=-\tfrac12,\qquad
\langle S_M,T'_M\rangle=(2+1)/\sqrt{12}=\tfrac{\sqrt3}{2}.
\]
The corresponding products with $S'_M$ are
$\langle T_M,S'_M\rangle=(2+1)/\sqrt{12}=\sqrt3/2$ and
$\langle S_M,S'_M\rangle=1/2$.
Using the second rows of the tables, in the basis $(D,E,F)$, gives
exactly the same four numbers at $M=-1/2$. Thus, for both magnetic
values,
\[
\boxed{\begin{pmatrix}T'_M\\S'_M\end{pmatrix}
=\begin{pmatrix}-1/2&\sqrt3/2\\ \sqrt3/2&1/2\end{pmatrix}
\begin{pmatrix}T_M\\S_M\end{pmatrix}}.
\]
The real matrix squares to the identity and has orthonormal rows.
It therefore changes basis in the same two-dimensional doublet space
at each $M$. Together with the unchanged quartet, it is a unitary
change of basis in the full eight-dimensional Hilbert space.
\erratum{The Instructor's Manual repeats the relative doublet-phase mismatch
from 16.2.3 in its triplet--electron coupling. The expansions here
use $J_-T_{1/2}=T_{-1/2}$ and $J_-T'_{1/2}=T'_{-1/2}$, so the
same recoupling matrix applies to both $M$ values.}

\Needspace{12\baselineskip}
\subsection*{16.3 Spherical Tensors}

\Needspace{5\baselineskip}
\par\bigskip
\noindent\textbf{16.3.1}\quad For $l=1$, Eq.~(16.52) becomes
\[
Y_{1}^{m}(\theta',\varphi')
=\sum_{m'=-1}^{1}D_{m'm}^{1}(\alpha,\beta,\gamma)\,
Y_{1}^{m'}(\theta,\varphi).
\]
Rewrite these spherical harmonics in Cartesian form. Show that the
resulting Cartesian coordinate equations are equivalent to the Euler
rotation matrix $\mathsf{A}(\alpha,\beta,\gamma)$, Eq.~(3.37).
\solution
Use the passive rotations of Eq.~(3.37): first $\alpha$ about $z$,
then $\beta$ about the intermediate $y$, then $\gamma$ about the
new $z$. Define
\[
Z(t)=\begin{pmatrix}\cos t&\sin t&0\\-\sin t&\cos t&0\\0&0&1\end{pmatrix},
\qquad
B(\beta)=\begin{pmatrix}\cos\beta&0&-\sin\beta\\0&1&0\\
\sin\beta&0&\cos\beta\end{pmatrix}.
\]
Thus the Cartesian column transforms by
$\mathbf r'=A\mathbf r$, where $A=Z(\gamma)B(\beta)Z(\alpha)$.

The $l=1$ harmonics are linear combinations of the Cartesian
components $(x,y,z)$ of $\hat{\mathbf r}$: since
$\sin\theta\,e^{\pm i\varphi}=x\pm iy$ and $\cos\theta=z$,
$Y_1^{\pm1}=\mp\sqrt{3/(8\pi)}\,(x\pm iy)$ and
$Y_1^0=\sqrt{3/(4\pi)}\,z$, that is,
\[
\mathcal Y=\begin{pmatrix}Y_1^{-1}\\Y_1^0\\Y_1^1\end{pmatrix}
=\sqrt{\frac3{4\pi}}\,Q\hat{\mathbf r},\qquad
Q=\begin{pmatrix}
1/\sqrt2&-i/\sqrt2&0\\
0&0&1\\
-1/\sqrt2&-i/\sqrt2&0
\end{pmatrix}.
\]
Direct row products give $QQ^\dagger=I$; for a square matrix this
also gives $Q^\dagger Q=I$, so
$\hat{\mathbf r}=\sqrt{4\pi/3}\,Q^\dagger\mathcal Y$.
In the index convention of Eq.~(16.52),
$Y_1^m(\Omega')=\sum_{m'}D^1_{m'm}Y_1^{m'}(\Omega)$, so the $m$th
component of $\mathcal Y'$ is $\sum_{m'}(D^1)^T_{mm'}\mathcal Y_{m'}$:
the matrix acting on this column is $(D^1)^T$. On the other hand,
\[
\mathcal Y'=\sqrt{\frac3{4\pi}}\,Q\hat{\mathbf r}'
=\sqrt{\frac3{4\pi}}\,QA\hat{\mathbf r}
=QAQ^\dagger\mathcal Y.
\]
The three functions $Y_1^m$ are linearly independent, so the two
matrices agree:
\[
\boxed{(D^1)^T=QAQ^\dagger,\quad A=Q^\dagger(D^1)^TQ}.
\]
The transpose here follows from the displayed summation indices.

To compare with the explicit rotation matrix $D^1$, compute the
elementary spherical-basis matrices
\[
QZ(t)Q^\dagger=E(t)=\operatorname{diag}(e^{it},1,e^{-it}),
\]
\[
QB(\beta)Q^\dagger=
H(\beta)=\begin{pmatrix}
(1+c_\beta)/2&-s_\beta/\sqrt2&(1-c_\beta)/2\\
s_\beta/\sqrt2&c_\beta&-s_\beta/\sqrt2\\
(1-c_\beta)/2&s_\beta/\sqrt2&(1+c_\beta)/2
\end{pmatrix},
\]
where $c_t=\cos t$, $s_t=\sin t$.
To obtain them, write $(u_{-1},u_0,u_1)^T=Q\hat{\mathbf r}$, i.e.
$u_{\mp1}=\pm(x\mp iy)/\sqrt2$, $u_0=z$; inversely
$x=(u_{-1}-u_1)/\sqrt2$, $iy=-(u_{-1}+u_1)/\sqrt2$, $z=u_0$.
Under $Z(t)$, $x'\mp iy'=(x\mp iy)(\cos t\pm i\sin t)
=e^{\pm it}(x\mp iy)$ and $z'=z$, which gives $E(t)$.
Under $B(\beta)$, $x'=c_\beta x-s_\beta z$, $y'=y$,
$z'=s_\beta x+c_\beta z$, so for example
\[
u'_{-1}=\frac{x'-iy'}{\sqrt2}
=\frac{c_\beta(u_{-1}-u_1)+(u_{-1}+u_1)}{2}-\frac{s_\beta u_0}{\sqrt2},
\]
which is the first row of $H$; the rows for $u'_0=z'$ and
$u'_1=-(x'+iy')/\sqrt2$ follow in the same way.
Inserting $Q^\dagger Q=I$ between the factors of
$QZ(\gamma)B(\beta)Z(\alpha)Q^\dagger$ gives
$(D^1)^T=E(\gamma)H(\beta)E(\alpha)$.
Since the diagonal entries of $E(t)$ are $e^{-imt}$ ($m=-1,0,1$),
this is, entry by entry, the standard Euler-angle form of the $l=1$
rotation matrix:
\[
D^1_{m'm}(\alpha,\beta,\gamma)
=e^{-im'\alpha}\,d^1_{m'm}(\beta)\,e^{-im\gamma},\qquad
d^1_{m'm}(\beta)=H_{mm'}(\beta),
\]
that is, $d^1_{11}=d^1_{-1,-1}=(1+c_\beta)/2$,
$d^1_{1,-1}=d^1_{-1,1}=(1-c_\beta)/2$, $d^1_{00}=c_\beta$,
$d^1_{10}=d^1_{0,-1}=-s_\beta/\sqrt2$ and
$d^1_{01}=d^1_{-1,0}=s_\beta/\sqrt2$.
Conversely, inserting this explicit $D^1$ in
$\mathcal Y'=(D^1)^T\mathcal Y$ and multiplying by
$\sqrt{4\pi/3}\,Q^\dagger$ gives the Cartesian equations
\[
\hat{\mathbf r}'=Q^\dagger(D^1)^TQ\,\hat{\mathbf r}
=Q^\dagger E(\gamma)Q\,Q^\dagger H(\beta)Q\,Q^\dagger E(\alpha)Q\,
\hat{\mathbf r}
=Z(\gamma)B(\beta)Z(\alpha)\,\hat{\mathbf r},
\]
because $Q^\dagger E(t)Q=Z(t)$ and $Q^\dagger H(\beta)Q=B(\beta)$ by
the unitarity of $Q$.
Multiplying out the two right-hand factors gives
\[
B(\beta)Z(\alpha)=
\begin{pmatrix}
c_\beta c_\alpha&c_\beta s_\alpha&-s_\beta\\
-s_\alpha&c_\alpha&0\\
s_\beta c_\alpha&s_\beta s_\alpha&c_\beta
\end{pmatrix}.
\]
Multiplying by $Z(\gamma)$ on the left replaces the first row by
$c_\gamma$ times the first row plus $s_\gamma$ times the second, and
the second row by $-s_\gamma$ times the first plus $c_\gamma$ times
the second; the third row is unchanged. This yields
\[
\boxed{A=
\begin{pmatrix}
c_\gamma c_\beta c_\alpha-s_\gamma s_\alpha&
c_\gamma c_\beta s_\alpha+s_\gamma c_\alpha&-c_\gamma s_\beta\\
-s_\gamma c_\beta c_\alpha-c_\gamma s_\alpha&
-s_\gamma c_\beta s_\alpha+c_\gamma c_\alpha&s_\gamma s_\beta\\
s_\beta c_\alpha&s_\beta s_\alpha&c_\beta
\end{pmatrix}}.
\]
This is exactly Eq.~(3.37), with the same passive-angle convention.

\Needspace{5\baselineskip}
\par\bigskip
\noindent\textbf{16.3.2}\quad In proving the addition theorem, we
assumed that $Y_{l}^{k}(\theta_{1},\varphi_{1})$ could be expanded in
a series of $Y_{l}^{m}(\theta_{2},\varphi_{2})$, in which $m$ varied
from $-l$ to $+l$ but $l$ was held fixed. What arguments can you
develop to justify summing only over the upper index, $m$, and not
over the lower index, $l$?
\par\emph{Hints.}~Examine the \erratum{homogeneous harmonic
polynomials $r^lY_l^m$}. Their restrictions to the sphere have terms
of the form $\cos^{l-p}\theta\sin^p\theta\erratum{e^{im\varphi}}$,
or $x^{l-p-s}y^pz^s/r^l$. Alternatively examine the rotation
invariance of the angular Laplacian. \erratum{The textbook's angular
hint omits the azimuthal factor; homogeneity alone also needs the
harmonicity condition to exclude lower-degree traces.}
\solution
An expansion of $Y_l^k(\theta_1,\varphi_1)$ in functions of
$(\theta_2,\varphi_2)$ requires $(\theta_1,\varphi_1)$ to be determined
by $(\theta_2,\varphi_2)$: the two pairs are the angles of one direction
referred to two sets of axes related by a fixed rotation. Let $R$ be
the rotation matrix with $\hat{\mathbf r}_1=R\hat{\mathbf r}_2$ for the
component columns, and let $U_Rf(\mathbf r)=f(R\mathbf r)$ denote its
action on functions; then $Y_l^k(\theta_1,\varphi_1)$, as a function
of $(\theta_2,\varphi_2)$, is $U_RY_l^k$. A rotation
preserves $r$, dot products, and the Cartesian Laplacian. The identity
\[
\nabla^2=\frac1{r^2}\frac{\partial}{\partial r}
\left(r^2\frac{\partial}{\partial r}\right)-\frac{\mathbf L^2}{r^2}
\]
can be solved as
$\mathbf L^2=\partial_r(r^2\partial_r)-r^2\nabla^2$. Each term on the
right commutes with $U_R$: $\nabla^2$ is rotation invariant, and
$r$ and $\partial_r$ are unchanged because $R(r\hat{\mathbf r})
=r\,R\hat{\mathbf r}$. Therefore $U_R\mathbf L^2=\mathbf L^2U_R$. Expand the
rotated harmonic in the complete orthonormal spherical-harmonic basis:
\[
U_RY_l^k=\sum_{n=0}^{\infty}\sum_{m=-n}^{n}a_{nm}Y_n^m.
\]
Applying $\mathbf L^2$ in the two orders, $\mathbf L^2U_RY_l^k
=U_R\mathbf L^2Y_l^k=l(l+1)U_RY_l^k$, gives
\[
l(l+1)\sum_{nm}a_{nm}Y_n^m
=\sum_{nm}n(n+1)a_{nm}Y_n^m.
\]
Multiply by $Y_{n'}^{m'*}$ and integrate, using orthonormality:
\[
\big[n'(n'+1)-l(l+1)\big]a_{n'm'}=0.
\]
(Equivalently, $n'(n'+1)a_{n'm'}=\langle\mathbf L^2Y_{n'}^{m'},U_RY_l^k\rangle
=\langle Y_{n'}^{m'},\mathbf L^2U_RY_l^k\rangle=l(l+1)a_{n'm'}$ by the
Hermiticity of $\mathbf L^2$, without differentiating the series.)
Since $n(n+1)$ increases strictly for $n\ge0$, the factor is zero only
for $n'=l$.
All other coefficients vanish, proving
\[
\boxed{Y_l^k(\theta_1,\varphi_1)
=\sum_{m=-l}^{l}a_{lm}Y_l^m(\theta_2,\varphi_2)}.
\]

The polynomial hint gives the same conclusion when its harmonic
condition is retained. The solid harmonic $r^lY_l^k$ is a homogeneous
harmonic polynomial of degree $l$. Linear substitution by a rotation
preserves its degree, and rotation invariance of $\nabla^2$ preserves
its being harmonic. It therefore remains in the space spanned by
$r^lY_l^{-l},\ldots,r^lY_l^l$: writing a homogeneous polynomial of
degree $l$ as $r^lf(\hat{\mathbf r})$, the displayed form of
$\nabla^2$ gives $\nabla^2(r^lf)=r^{l-2}[l(l+1)f-\mathbf L^2f]$, so
harmonicity means $\mathbf L^2f=l(l+1)f$, and the coefficient argument
above leaves only the $Y_l^m$. Homogeneity alone would not suffice:
a general homogeneous polynomial may contain traces such as
$r^2$ times a lower-degree harmonic. Harmonicity excludes those traces.

\Needspace{5\baselineskip}
\par\bigskip
\noindent\textbf{16.3.3}\quad An atomic electron with angular
momentum $l$ and magnetic quantum number $m$ has a wave function
\[
\psi(r,\theta,\varphi)=f(r)Y_{l}^{m}(\theta,\varphi).
\]
Show that the sum of the electron densities in a given complete shell
is spherically symmetric; that is,
$\sum_{m=-l}^{l}\psi^{*}(r,\theta,\varphi)\psi(r,\theta,\varphi)$ is
independent of $\theta$ and $\varphi$.
\solution
All members of the specified shell have the same radial factor.
Consequently
\[
\sum_{m=-l}^{l}\psi_m^*\psi_m
=|f(r)|^2\sum_{m=-l}^{l}|Y_l^m(\theta,\varphi)|^2.
\]
Apply the addition theorem to two identical directions. Their
separation angle is zero, so $\cos\gamma=1$, and $P_l(1)=1$:
\[
1=P_l(1)=\frac{4\pi}{2l+1}
\sum_{m=-l}^{l}Y_l^{m*}(\theta,\varphi)Y_l^m(\theta,\varphi),
\]
that is, $\sum_{m}|Y_l^m(\theta,\varphi)|^2=(2l+1)/(4\pi)$ for every
direction. Thus
\[
\boxed{\sum_{m=-l}^{l}\psi_m^*\psi_m
=\frac{2l+1}{4\pi}|f(r)|^2}.
\]
This depends only on $r$. Equal occupancy of all magnetic states
is the complete-shell condition used by the stated sum; a common
spin-occupancy factor merely multiplies this radial density.

\Needspace{5\baselineskip}
\par\bigskip
\noindent\textbf{16.3.4}\quad The \erratum{potential energy} of an electron at point
$\mathbf{r}_{e}$ in the field of $Z$ protons at points
$\mathbf{r}_{p}$ is
\[
\Phi=-\frac{e^{2}}{4\pi\varepsilon_{0}}
\sum_{p=1}^{Z}\frac{1}{|\mathbf{r}_{e}-\mathbf{r}_{p}|}.
\]
Show that for $r_{e}$ larger than all $r_{p}$, this may be written as
\[
\Phi=-\frac{e^{2}}{4\pi\varepsilon_{0}r_{e}}
\sum_{p=1}^{Z}\sum_{L,M}\left(\frac{r_{p}}{r_{e}}\right)^{L}
\frac{4\pi}{2L+1}\,
Y_{L}^{M}(\theta_{p},\varphi_{p})^{*}Y_{L}^{M}(\theta_{e},\varphi_{e}).
\]
How should $\Phi$ be written for $r_{e}<r_{p}$?
\solution
\erratum{The displayed expression contains the charge product
$(-e)(+e)=-e^2$, so it is an interaction energy. The electric
potential of the protons alone has coefficient $+e/(4\pi\varepsilon_0)$.}
For each proton define $\gamma_p$ as the angle between
$\mathbf r_e$ and $\mathbf r_p$, so that
$|\mathbf r_e-\mathbf r_p|^2=r_e^2-2r_er_p\cos\gamma_p+r_p^2$.
With $t=r_p/r_e<1$, the Legendre generating function gives
\[
\frac1{|\mathbf r_e-\mathbf r_p|}
=\frac1{r_e}(1-2t\cos\gamma_p+t^2)^{-1/2}
=\frac1{r_e}\sum_{L=0}^{\infty}t^LP_L(\cos\gamma_p).
\]
The addition theorem is
\[
P_L(\cos\gamma_p)=\frac{4\pi}{2L+1}
\sum_{M=-L}^{L}Y_L^{M*}(\Omega_p)Y_L^M(\Omega_e).
\]
Substitution into each term of the proton sum, with the factor
$-e^2/(4\pi\varepsilon_0)$, gives exactly the first requested
expansion.

For $r_e<r_p$, factor out $r_p$ instead:
$|\mathbf r_e-\mathbf r_p|^{-1}
=r_p^{-1}(1-2t'\cos\gamma_p+t'^2)^{-1/2}$ with $t'=r_e/r_p<1$.
The same angular factor then gives
\[
\boxed{\Phi=-\frac{e^2}{4\pi\varepsilon_0}
\sum_{p=1}^{Z}\frac1{r_p}
\sum_{L=0}^{\infty}\left(\frac{r_e}{r_p}\right)^L
\frac{4\pi}{2L+1}\sum_{M=-L}^{L}
Y_L^{M*}(\Omega_p)Y_L^M(\Omega_e)}
\]
when the electron lies inside every proton radius. More generally,
the expansion can be used term by term, even when some protons
are inside $r_e$ and others outside:
\[
\boxed{\Phi=-\frac{e^2}{4\pi\varepsilon_0}
\sum_{p=1}^{Z}\sum_{L=0}^{\infty}
\frac{r_{<p}^{\,L}}{r_{>p}^{\,L+1}}\frac{4\pi}{2L+1}
\sum_{M=-L}^{L}Y_L^{M*}(\Omega_p)Y_L^M(\Omega_e)},
\]
where $r_{<p}=\min(r_e,r_p)$ and $r_{>p}=\max(r_e,r_p)$.
The strictly unequal radii assumed above make each generating
series convergent. Reversing which radius is larger changes only
the radial factor, not the angular addition theorem.

\Needspace{5\baselineskip}
\par\bigskip
\noindent\textbf{16.3.5}\quad Two protons are uniformly distributed
within the same spherical volume. If the coordinates of one element
of charge are $(r_{1},\theta_{1},\varphi_{1})$ and the coordinates of
the other are $(r_{2},\theta_{2},\varphi_{2})$ and $r_{12}$ is the
distance between them, the element of repulsion energy will be given
by
\[
d\psi=\rho^{2}\frac{d\tau_{1}\,d\tau_{2}}{r_{12}}
=\rho^{2}\frac{r_{1}^{2}dr_{1}\sin\theta_{1}\,d\theta_{1}\,d\varphi_{1}\,
r_{2}^{2}dr_{2}\sin\theta_{2}\,d\theta_{2}\,d\varphi_{2}}{r_{12}},
\]
where
\[
\rho=\frac{\text{charge}}{\text{volume}}=\frac{3e}{4\pi R^{3}}
\qquad\text{and}\qquad
r_{12}^{2}=r_{1}^{2}+r_{2}^{2}-2r_{1}r_{2}\cos\gamma.
\]
Here $\rho$ is the charge density and $\gamma$ is the angle between
$\mathbf{r}_{1}$ and $\mathbf{r}_{2}$. Calculate the total
electrostatic energy (of repulsion) of the two protons. This
calculation is used in accounting for the mass difference in
``mirror'' nuclei, such as O${}^{15}$ and N${}^{15}$.
\par\nobreak\hfill\emph{ANS.}\enspace$\dfrac{6}{5}\dfrac{e^{2}}{R}$.
\solution
The angular average of the Coulomb expansion retains only $L=0$.
As in 16.3.4, for $r_1\ne r_2$,
\[
\frac1{r_{12}}=\sum_{L=0}^{\infty}\frac{r_<^L}{r_>^{L+1}}
P_L(\cos\gamma),\qquad r_<=\min(r_1,r_2),\ r_>=\max(r_1,r_2).
\]
By the addition theorem and $\int Y_L^M\,d\Omega
=\sqrt{4\pi}\int Y_0^{0*}Y_L^M\,d\Omega=\sqrt{4\pi}\,\delta_{L0}\delta_{M0}$,
\[
\int P_L(\hat{\mathbf r}_1\cdot\hat{\mathbf r}_2)\,d\Omega_2
=\frac{4\pi}{2L+1}\sum_{M}Y_L^{M*}(\Omega_1)\sqrt{4\pi}\,
\delta_{L0}\delta_{M0}=4\pi\delta_{L0},
\]
since $Y_0^0=1/\sqrt{4\pi}$. The remaining $\Omega_1$ integral
gives another factor $4\pi$, so
\[
\int d\Omega_1\int d\Omega_2\,\frac1{r_{12}}
=\frac{16\pi^2}{\max(r_1,r_2)}.
\]
After these angular integrations, $U=\int d\psi$ is the double radial
integral in the first line below. The diagonal $r_1=r_2$ has zero
measure in it. The integrand is symmetric in $r_1,r_2$, and on the
triangle $r_2<r_1$ we have $\max(r_1,r_2)=r_1$. Split its square
into the two equal triangles:
\[
\begin{aligned}
U&=16\pi^2\rho^2\int_0^R\!\!\int_0^R
\frac{r_1^2r_2^2}{\max(r_1,r_2)}\,dr_2\,dr_1\\
&=32\pi^2\rho^2\int_0^R r_1\,dr_1
\int_0^{r_1}r_2^2\,dr_2\\
&=32\pi^2\rho^2\int_0^R r_1\,\frac{r_1^3}{3}\,dr_1
=\frac{32\pi^2\rho^2}{3}\int_0^R r_1^4\,dr_1
=\frac{32\pi^2\rho^2R^5}{15}.
\end{aligned}
\]
Substitute $\rho^2=9e^2/(16\pi^2R^6)$:
\[
\boxed{U=\frac{32\cdot9}{15\cdot16}\frac{e^2}{R}
=\frac{288}{240}\frac{e^2}{R}
=\frac65\frac{e^2}{R}}.
\]
There is no factor $1/2$: this is the mutual energy of two distinct
charge distributions, and each pair with one element from each
distribution is counted once. The problem uses a Coulomb constant
of one; in SI units the result is multiplied by $1/(4\pi\varepsilon_0)$.

\Needspace{5\baselineskip}
\par\bigskip
\noindent\textbf{16.3.6}\quad Each of the two $1s$ electrons in
helium may be described by a hydrogenic wave function
\[
\psi(r)=\left(\frac{Z^{3}}{\pi a_{0}^{3}}\right)^{1/2}e^{-Zr/a_{0}}
\]
in the absence of the other electron. Here $Z$, the atomic number, is
2. The symbol $a_{0}$ is the Bohr radius, $\hbar^{2}/me^{2}$. Find
the mutual potential energy of the two electrons, given by
\[
\int\psi^{*}(\mathbf{r}_{1})\psi^{*}(\mathbf{r}_{2})
\frac{e^{2}}{|\mathbf{r}_{1}-\mathbf{r}_{2}|}\,
\psi(\mathbf{r}_{1})\psi(\mathbf{r}_{2})\,d^{3}r_{1}\,d^{3}r_{2}.
\]
\par\nobreak\hfill\emph{ANS.}\enspace$\dfrac{5e^{2}Z}{8a_{0}}$.
\solution
Put $\lambda=2Z/a_0$ and $N=Z^3/(\pi a_0^3)$, so
$|\psi(r)|^2=Ne^{-\lambda r}$. Its normalization is
\[
4\pi N\int_0^\infty r^2e^{-\lambda r}\,dr
=4\pi N\,\frac{2}{\lambda^3}
=\frac{8\pi Z^3}{\pi a_0^3}\,\frac{a_0^3}{8Z^3}=1.
\]
The angular integration is the same as in 16.3.5: it replaces
$1/|\mathbf r_1-\mathbf r_2|$ by $16\pi^2/\max(r_1,r_2)$, so
\[
U=16\pi^2e^2N^2\int_0^\infty\!\!\int_0^\infty
\frac{r_1^2r_2^2e^{-\lambda(r_1+r_2)}}{\max(r_1,r_2)}\,dr_2\,dr_1.
\]
The integrand is symmetric, so splitting the radial quadrant into
the equal halves $r_2<r_1$ and $r_1<r_2$ gives
\[
U=32\pi^2e^2N^2
\int_0^\infty r_1e^{-\lambda r_1}\,dr_1
\int_0^{r_1}r_2^2e^{-\lambda r_2}\,dr_2.
\]
Set $x=\lambda r_2$, $y=\lambda r_1$, so that
$r_1\,dr_1\,r_2^2\,dr_2=\lambda^{-5}y\,dy\,x^2\,dx$ and the region
$0<r_2<r_1$ becomes $0<x<y$, and reverse the order:
\[
U=\frac{32\pi^2e^2N^2}{\lambda^5}
\int_0^\infty x^2e^{-x}
\left(\int_x^\infty ye^{-y}\,dy\right)dx.
\]
Integration by parts gives
\[
\int_x^\infty ye^{-y}dy
=[-ye^{-y}]_x^\infty+\int_x^\infty e^{-y}dy
=(x+1)e^{-x}.
\]
The $x$ integrand is therefore $x^2(x+1)e^{-2x}=(x^3+x^2)e^{-2x}$.
For $I_k=\int_0^\infty x^ke^{-2x}dx$ and $k\ge1$,
\[
I_k=\left[-\frac{x^ke^{-2x}}2\right]_0^\infty
+\frac{k}{2}\int_0^\infty x^{k-1}e^{-2x}dx
=\frac{k}{2}I_{k-1}.
\]
Since $I_0=1/2$, this gives $I_1=1/4$, $I_2=1/4$, and
$I_3=3/8$, or generally $I_k=k!/2^{k+1}$.
Hence the dimensionless integral is
\[
\int_0^\infty(x^3+x^2)e^{-2x}\,dx
=\frac{3!}{2^4}+\frac{2!}{2^3}=\frac38+\frac14=\frac58.
\]
The prefactor simplifies without an approximation, using
$N^2=Z^6/(\pi^2a_0^6)$ and $(2Z/a_0)^5=32Z^5/a_0^5$:
\[
\frac{32\pi^2e^2}{(2Z/a_0)^5}
\frac{Z^6}{\pi^2a_0^6}
=\frac{32e^2a_0^5}{32Z^5}\,\frac{Z^6}{a_0^6}
=\frac{e^2Z}{a_0}.
\]
Thus
\[
\boxed{U=\frac{5e^2Z}{8a_0}
=\frac{5e^2}{4a_0}\quad(Z=2)}.
\]
This evaluates the stated independent hydrogenic-orbital integral;
it uses the same Coulomb-unit convention as the problem.

\Needspace{5\baselineskip}
\par\bigskip
\noindent\textbf{16.3.7}\quad The probability of finding a $1s$
hydrogen electron in a volume element
$r^{2}dr\sin\theta\,d\theta\,d\varphi$ is
\[
\frac{1}{\pi a_{0}^{3}}\,e^{-2r/a_{0}}\,
r^{2}dr\sin\theta\,d\theta\,d\varphi,
\]
where $r$ is the distance of the electron from the nucleus. Find the
electrostatic potential of this charge distribution at points
$\mathbf{r}_{1}$, where you may not assume that $\mathbf{r}_{1}$ is
on the polar axis of your coordinate system. Calculate the potential
from
\[
V(\mathbf{r}_{1})=-\frac{e}{4\pi\varepsilon_{0}}
\int\frac{\rho(\mathbf{r}_{2})}{r_{12}}\,d^{3}r_{2},
\]
where $r_{12}=|\mathbf{r}_{1}-\mathbf{r}_{2}|$ and
$\rho(r)=e^{-2r/a_0}/(\pi a_0^3)$ is the probability density.
Expand \erratum{$1/r_{12}$}.
Apply the Legendre polynomial addition theorem and show that the
angular dependence of $V(\mathbf{r}_{1})$ drops out.
\par\nobreak\hfill\emph{ANS.}\enspace
$V(r_{1})=-\dfrac{e}{4\pi\varepsilon_{0}}
\left[\dfrac{1}{2r_{1}}\gamma\left(3,\dfrac{2r_{1}}{a_{0}}\right)
+\dfrac{1}{a_{0}}\Gamma\left(2,\dfrac{2r_{1}}{a_{0}}\right)\right]$, where
$\gamma$ and $\Gamma$ are incomplete gamma functions, Eq.~(13.73).
\solution
\erratum{The textbook's instruction ``Expand $r_{12}$'' must instead say
``Expand $1/r_{12}$'', since $1/r_{12}$ is the Coulomb kernel in the
integral.}
Write $R=r_1$ and $r=r_2$. Expand the inverse separation:
\[
\frac1{r_{12}}=\sum_{l=0}^{\infty}\frac{r_<^l}{r_>^{l+1}}
\frac{4\pi}{2l+1}\sum_{m=-l}^{l}
Y_l^m(\Omega_1)Y_l^{m*}(\Omega_2).
\]
The density is independent of $\Omega_2$. Because
$\int Y_l^{m*}\,d\Omega_2=\sqrt{4\pi}\delta_{l0}\delta_{m0}$
and $Y_0^0(\Omega_1)=1/\sqrt{4\pi}$, its angular integral is
$4\pi\cdot Y_0^0(\Omega_1)\sqrt{4\pi}/r_>=4\pi/r_>$. Thus no choice of
polar axis is needed. With $d^3r_2=r^2dr\,d\Omega_2$, the factor
$4\pi/(\pi a_0^3)=4/a_0^3$, and $r_>=R$ for $r<R$, $r_>=r$ for $r>R$,
\[
V(R)=-\frac{e}{4\pi\varepsilon_0}\frac4{a_0^3}
\left[\frac1R\int_0^R r^2e^{-2r/a_0}\,dr
+\int_R^\infty re^{-2r/a_0}\,dr\right].
\]
Put $q=2R/a_0$ and $r=a_0t/2$, so $dr=(a_0/2)\,dt$. By the definitions
$\gamma(s,q)=\int_0^q t^{s-1}e^{-t}dt$ and
$\Gamma(s,q)=\int_q^\infty t^{s-1}e^{-t}dt$,
\[
\int_0^R r^2e^{-2r/a_0}\,dr=\frac{a_0^3}{8}\,\gamma(3,q),\qquad
\int_R^\infty re^{-2r/a_0}\,dr=\frac{a_0^2}{4}\,\Gamma(2,q),
\]
and multiplying by $4/a_0^3$ gives
\[
\boxed{V(R)=-\frac{e}{4\pi\varepsilon_0}
\left[\frac{\gamma(3,q)}{2R}+\frac{\Gamma(2,q)}{a_0}\right]}.
\]
For clarity, the required integrations by parts are
\[
\begin{aligned}
\int_0^q te^{-t}dt
&=[-te^{-t}]_0^q+\int_0^q e^{-t}dt
=1-(q+1)e^{-q},\\
\gamma(3,q)=\int_0^q t^2e^{-t}dt
&=[-t^2e^{-t}]_0^q+2\int_0^q te^{-t}dt\\
&=2-e^{-q}(q^2+2q+2),\\
\Gamma(2,q)=\int_q^\infty te^{-t}dt
&=[-te^{-t}]_q^\infty+\int_q^\infty e^{-t}dt
=e^{-q}(q+1).
\end{aligned}
\]
Since $q/(2R)=1/a_0$, so that $q^2/(2R)=q/a_0$ and $2q/(2R)=2/a_0$,
substitution gives
\[
\frac{\gamma(3,q)}{2R}+\frac{\Gamma(2,q)}{a_0}
=\frac1R-e^{-q}\left(\frac{q}{a_0}+\frac2{a_0}+\frac1R\right)
+e^{-q}\left(\frac{q}{a_0}+\frac1{a_0}\right),
\]
in which the $q/a_0$ terms cancel and $2/a_0-1/a_0=1/a_0$, leaving
\[
\boxed{V(R)=-\frac{e}{4\pi\varepsilon_0}
\left[\frac1R-e^{-2R/a_0}\left(\frac1R+\frac1{a_0}\right)\right]}.
\]
For example, writing the bracket as
$(1-e^{-2R/a_0})/R-e^{-2R/a_0}/a_0$ and using
$(1-e^{-2R/a_0})/R\to2/a_0$ as $R\to0$ (since $1-e^{-x}=x+O(x^2)$)
gives the limit
$2/a_0-1/a_0$, so $V(0)=-e/(4\pi\varepsilon_0a_0)$. At large $R$ the
exponential disappears and the potential is that of charge $-e$,
as required.

\Needspace{5\baselineskip}
\par\bigskip
\noindent\textbf{16.3.8}\quad A hydrogen electron in a $2p$ orbital
has a charge distribution
\[
\rho=-\frac{e}{64\pi a_{0}^{5}}\,r^{2}e^{-r/a_{0}}\sin^{2}\theta,
\]
where $a_{0}=\hbar^{2}/me^{2}$ is the Bohr radius, and $r$ is the
distance between the electron and the nucleus. Find the electrostatic
\erratum{electron--nucleus interaction energy, treating the nucleus
as a point charge $+e$ at the origin. Also find the electric potential
of the electron distribution at an arbitrary field point.}
\solution
\erratum{For the hydrogen atom, the stated electrostatic energy is
interpreted as the electron--nucleus interaction. The manual instead
displays a position-dependent electric potential; both quantities
are obtained here. A classical self-energy
$\tfrac12\int\rho V\,d^3r$ would be a different quantity, so a
second charge is not required for every notion of electrostatic energy.}

Let the field point be $(R,\Theta,\Phi)$ and write $a=a_0$.
The identity $\sin^2\theta=\frac23[1-P_2(\cos\theta)]$ gives
\[
\rho(r,\theta)=-\frac{e}{96\pi a^5}r^2e^{-r/a}
[1-P_2(\cos\theta)].
\]
Its total charge is
\[
-\frac{e}{96\pi a^5}\,(4\pi)\int_0^\infty r^4e^{-r/a}dr
=-\frac{e}{24a^5}(4!a^5)=-e.
\]
Only the monopole and quadrupole survive. The addition theorem and
orthogonality give, for $l=0,2$,
\[
\int\frac{P_l(\cos\theta)}{|\mathbf R-\mathbf r|}\,d\Omega
=\frac{4\pi}{2l+1}\frac{r_<^l}{r_>^{l+1}}P_l(\cos\Theta).
\]
To see the factor, write
$P_l(\cos\theta)=\sqrt{4\pi/(2l+1)}Y_l^0(\Omega)$ in the
Coulomb expansion: the angular integral selects its $l,0$ term.
Consequently
\[
V(R,\Theta)=-\frac{e}{4\pi\varepsilon_0}\frac1{24a^5}
\left[I_0(R)-\frac{P_2(\cos\Theta)}5I_2(R)\right],
\]
where splitting each radial integral at $r=R$ gives
\[
\begin{aligned}
I_0(R)&=\frac1R\int_0^R r^4e^{-r/a}dr
+\int_R^\infty r^3e^{-r/a}dr,\\
I_2(R)&=\frac1{R^3}\int_0^R r^6e^{-r/a}dr
+R^2\int_R^\infty re^{-r/a}dr.
\end{aligned}
\]
Thus, with $q=R/a$,
\[
\boxed{\begin{aligned}
V(R,\Theta)=-\frac{e}{4\pi\varepsilon_0}
\bigg\{&
\frac1{24}\left[\frac{\gamma(5,q)}R+
\frac{\Gamma(4,q)}a\right]\\
&-\frac{P_2(\cos\Theta)}{120}
\left[\frac{a^2\gamma(7,q)}{R^3}
+\frac{R^2\Gamma(2,q)}{a^3}\right]\bigg\}.
\end{aligned}}
\]
\erratum{The manual's potential on printed p.~265 lacks the
overall factor $-e$ required by the given electron charge density.
The displayed result restores both the charge and its negative sign.}
A test charge $q_{\rm test}$ has energy $q_{\rm test}V(R,\Theta)$.

For the hydrogen nucleus, $q_{\rm test}=+e$ and $R=0$.
Since $e^{-r/a}=1+O(r)$ near the origin, the interior contributions are
\[
\frac1R\int_0^R r^4e^{-r/a}dr
=\frac{R^4}{5}+O(R^5),\qquad
\frac1{R^3}\int_0^R r^6e^{-r/a}dr
=\frac{R^4}{7}+O(R^5).
\]
Both vanish; the exterior quadrupole contribution is
$R^2\int_R^\infty re^{-r/a}dr=O(R^2)$ and also vanishes.
The surviving contribution in the braces is
$\Gamma(4,0)/(24a)=3!/(24a)=1/(4a)$. Therefore
\[
\boxed{U_{\text{electron--nucleus}}
=eV(0)=-\frac{e^2}{4\pi\varepsilon_0}\frac1{4a_0}}.
\]
Equivalently one may integrate the nuclear potential directly:
\[
U=\int\rho\,\frac{e}{4\pi\varepsilon_0 r}\,d^3r
=-\frac{e^2}{4\pi\varepsilon_0}\frac1{24a^5}
\int_0^\infty r^3e^{-r/a}dr
=-\frac{e^2}{4\pi\varepsilon_0}\frac1{4a}.
\]
In the Coulomb units used in the question's $a_0=\hbar^2/(me^2)$,
set $1/(4\pi\varepsilon_0)=1$: the energy is $-e^2/(4a_0)$.
This is the electrostatic interaction energy; the kinetic energy
is not part of the stated electrostatic calculation.

\Needspace{5\baselineskip}
\par\bigskip
\noindent\textbf{16.3.9}\quad
\erratum{Interpret the delta and both sums as distributions with
respect to $d\Omega=\sin\theta\,d\theta\,d\varphi$.}
\begin{enumerate}
\item[(a)] As a Laplace series and as an example of Eq.~(5.27), show
that
\[
\delta(\Omega_{1}-\Omega_{2})
=\sum_{l=0}^{\infty}\sum_{m=-l}^{l}
Y_{l}^{m}(\theta_{2},\varphi_{2})^{*}Y_{l}^{m}(\theta_{1},\varphi_{1}).
\]
\item[(b)] Show also that this same representation of the Dirac delta
function may be written as
\[
\delta(\Omega_{1}-\Omega_{2})
=\sum_{l=0}^{\infty}\frac{2l+1}{4\pi}\,P_{l}(\cos\gamma),
\]
and identify $\gamma$. Now, if you can justify equating the summations
over $l$ term by term, you have an alternate derivation of the
spherical harmonic addition theorem.
\end{enumerate}
\solution
\erratum{These infinite sums represent a distribution on the sphere,
with measure $d\Omega=\sin\theta\,d\theta\,d\varphi$; they are not
ordinary pointwise convergent formulas for a function.}

\textbf{(a)} For a smooth test function $f$, let
$f_{lm}=\int Y_l^{m*}(\Omega_2)f(\Omega_2)\,d\Omega_2$.
Completeness gives
$f(\Omega_1)=\sum_{lm}f_{lm}Y_l^m(\Omega_1)$.
Therefore
\[
\int\left[\sum_{lm}Y_l^{m*}(\Omega_2)Y_l^m(\Omega_1)\right]
f(\Omega_2)\,d\Omega_2
=\sum_{lm}f_{lm}Y_l^m(\Omega_1)=f(\Omega_1).
\]
This is precisely the defining action of
$\delta(\Omega_1-\Omega_2)$, proving (a) in the distributional sense.

\textbf{(b)} Define the finite-dimensional projection kernel
\[
K_l(\Omega_1,\Omega_2)=
\sum_{m=-l}^{l}Y_l^m(\Omega_1)Y_l^{m*}(\Omega_2).
\]
A common rotation of the two directions mixes their $m$ components
by the same unitary matrix: rotations preserve $d\Omega$ and hence
the harmonic inner product. In the sum, the two matrices multiply
to the identity. Thus $K_l$ is unchanged by a common rotation.
Choose one taking $\Omega_2$ to the north pole. The other direction
then has polar angle $\gamma$, their angular separation. At the pole,
\[
Y_l^m(0,\varphi)=0\quad(m\ne0),\qquad
Y_l^0(0,\varphi)=\sqrt{\frac{2l+1}{4\pi}}.
\]
It follows directly, without assuming the addition theorem, that
\[
K_l(\Omega_1,\Omega_2)
=Y_l^0(\gamma,0)\sqrt{\frac{2l+1}{4\pi}}
=\boxed{\frac{2l+1}{4\pi}P_l(\cos\gamma)},
\]
where
\[
\cos\gamma=\cos\theta_1\cos\theta_2+
\sin\theta_1\sin\theta_2\cos(\varphi_1-\varphi_2).
\]
Summing these kernels gives the stated Legendre representation
of the delta distribution.

The term-by-term identification is justified by orthogonal
projection, rather than by subtracting two arbitrary infinite sums.
Apply the degree-$l$ projector in the first angular variable.
It annihilates all $K_k$ with $k\ne l$, since
$\int Y_l^{m*}Y_k^q\,d\Omega=\delta_{lk}\delta_{mq}$, and leaves
$K_l$ unchanged. Thus each degree block is uniquely determined.
The boxed identity is the addition theorem.
\erratum{The manual (printed p.~266) says that reconstructing the
arbitrary test function yields $Y_l^m(\Omega_1)$. Its preceding
sum instead yields $f(\Omega_1)$, as the delta's defining action
requires. A single harmonic is recovered only when $f$ is that harmonic.}

\Needspace{5\baselineskip}
\par\bigskip
\noindent\textbf{16.3.10}\quad Verify the following, with
\erratum{$L\ge0$ and $|M|\le L$; in (d) also require
$L\ge1$ and $|M+1|\le L-1$}:
\begin{enumerate}
\item[(a)] $\displaystyle\int
Y_{L}^{M}(\theta,\varphi)Y_{0}^{0}(\theta,\varphi)
Y_{L}^{M*}(\theta,\varphi)\,d\Omega=\frac{1}{\sqrt{4\pi}}$,
\item[(b)] $\displaystyle\int Y_{L}^{M}Y_{1}^{0}Y_{L+1}^{M*}\,d\Omega
=\sqrt{\frac{3}{4\pi}\,
\frac{(L+M+1)(L-M+1)}{(2L+1)(2L+3)}}$,
\item[(c)] $\displaystyle\int
Y_{L}^{M}Y_{1}^{1}Y_{L+1}^{M+1*}\,d\Omega
=\sqrt{\frac{3}{8\pi}\,
\frac{(L+M+1)(L+M+2)}{(2L+1)(2L+3)}}$,
\item[(d)] $\displaystyle\int
Y_{L}^{M}Y_{1}^{1}Y_{L-1}^{M+1*}\,d\Omega
=-\sqrt{\frac{3}{8\pi}\,
\frac{(L-M)(L-M-1)}{(2L-1)(2L+1)}}$.
\end{enumerate}
These integrals were used in an investigation of the angular
correlation of internal conversion electrons.
\solution
Let
\[
N_{LM}=\sqrt{\frac{2L+1}{4\pi}\frac{(L-M)!}{(L+M)!}},
\qquad Y_L^M=N_{LM}P_L^M(\cos\theta)e^{iM\varphi}.
\]
All integrals use normalized harmonics and $d\Omega$.
\erratum{Each harmonic appearing in an integral must have
nonnegative degree and $|m|\le l$. In particular, (d) is an
integral of the displayed harmonics only when $L\ge1$ and
$|M+1|\le L-1$; absent terms in recurrence formulas are omitted.}

\textbf{(a)} Since $Y_0^0=1/\sqrt{4\pi}$,
\[
\int Y_L^MY_0^0Y_L^{M*}d\Omega
=\frac1{\sqrt{4\pi}}\int|Y_L^M|^2d\Omega
=\boxed{\frac1{\sqrt{4\pi}}}.
\]

\textbf{(b)} Put $x=\cos\theta$. Equation (15.88) gives
\[
xP_L^M=\frac{L-M+1}{2L+1}P_{L+1}^M+
\frac{L+M}{2L+1}P_{L-1}^M.
\]
For the upper term, the normalization ratio is
\[
\frac{N_{LM}}{N_{L+1,M}}
=\sqrt{\frac{2L+1}{2L+3}\frac{L+M+1}{L-M+1}}.
\]
Multiplying by $(L-M+1)/(2L+1)$ therefore gives
\[
A_+=\sqrt{\frac{(L+1)^2-M^2}{(2L+1)(2L+3)}}.
\]
For the lower term,
\[
\frac{N_{LM}}{N_{L-1,M}}
=\sqrt{\frac{2L+1}{2L-1}\frac{L-M}{L+M}},
\qquad
A_-=\sqrt{\frac{L^2-M^2}{(2L-1)(2L+1)}}.
\]
At an endpoint the zero coefficient removes the nonexistent term.
Thus
$\cos\theta\,Y_L^M=A_+Y_{L+1}^M+A_-Y_{L-1}^M$.
Multiply by $Y_{L+1}^{M*}$, integrate, and use
$Y_1^0=\sqrt{3/(4\pi)}\cos\theta$:
\[
\boxed{\int Y_L^MY_1^0Y_{L+1}^{M*}d\Omega
=\sqrt{\frac3{4\pi}
\frac{(L+M+1)(L-M+1)}{(2L+1)(2L+3)}}}.
\]

\textbf{(c),(d)} Equation (15.89) is
\[
\sqrt{1-x^2}\,P_L^M=
\frac{P_{L-1}^{M+1}-P_{L+1}^{M+1}}{2L+1}.
\]
The required ratios, including the azimuthal factor $e^{i\varphi}$,
are
\[
\begin{aligned}
\frac{N_{LM}}{(2L+1)N_{L+1,M+1}}
&=\sqrt{\frac{(L+M+1)(L+M+2)}{(2L+1)(2L+3)}}=B_+,\\
\frac{N_{LM}}{(2L+1)N_{L-1,M+1}}
&=\sqrt{\frac{(L-M)(L-M-1)}{(2L-1)(2L+1)}}=B_-.
\end{aligned}
\]
For example the upper factorial quotient cancels $(L-M)!$
and leaves $(L+M+2)!/(L+M)!=(L+M+2)(L+M+1)$;
the lower quotient similarly leaves $(L-M)(L-M-1)$.
Hence
\[
\sin\theta\,e^{i\varphi}Y_L^M
=-B_+Y_{L+1}^{M+1}+B_-Y_{L-1}^{M+1}.
\]
Because $Y_1^1=-\sqrt{3/(8\pi)}\sin\theta\,e^{i\varphi}$,
orthogonality now selects each coefficient:
\[
\boxed{\begin{aligned}
\int Y_L^MY_1^1Y_{L+1}^{M+1*}d\Omega
&=\sqrt{\frac3{8\pi}
\frac{(L+M+1)(L+M+2)}{(2L+1)(2L+3)}},\\
\int Y_L^MY_1^1Y_{L-1}^{M+1*}d\Omega
&=-\sqrt{\frac3{8\pi}
\frac{(L-M)(L-M-1)}{(2L-1)(2L+1)}}.
\end{aligned}}
\]

\Needspace{5\baselineskip}
\par\bigskip
\noindent\textbf{16.3.11}\quad For \erratum{nonnegative integers
$L,N$}, show that
\begin{enumerate}
\item[(a)] $\displaystyle\int_{-1}^{1}xP_{L}(x)P_{N}(x)\,dx
=\begin{cases}
\dfrac{2(L+1)}{(2L+1)(2L+3)}, & N=L+1,\\[12pt]
\dfrac{2L}{(2L-1)(2L+1)}, & N=L-1,\\[8pt]
\erratum{0}, & \erratum{\text{all other }N\ge0}.
\end{cases}$
\item[(b)] $\displaystyle\int_{-1}^{1}x^{2}P_{L}(x)P_{N}(x)\,dx
=\begin{cases}
\dfrac{2(L+1)(L+2)}{(2L+1)(2L+3)(2L+5)}, & N=L+2,\\[12pt]
\dfrac{2(2L^{2}+2L-1)}{(2L-1)(2L+1)(2L+3)}, & N=L,\\[12pt]
\dfrac{2L(L-1)}{(2L-3)(2L-1)(2L+1)}, & N=L-2,\\[8pt]
\erratum{0}, & \erratum{\text{all other }N\ge0}.
\end{cases}$
\end{enumerate}
\solution
Take $L,N$ to be nonnegative integers. Write the three-term recurrence
as
\[
xP_L=a_LP_{L+1}+b_LP_{L-1},\qquad
a_L=\frac{L+1}{2L+1},\quad b_L=\frac{L}{2L+1}.
\]
At $L=0$ this is $xP_0=P_1$; the negative-degree term is omitted.
Orthogonality is
$\int_{-1}^1P_jP_N\,dx=2\delta_{jN}/(2N+1)$.

\textbf{(a)} Multiplying the recurrence by $P_N$ and integrating gives
\[
\int_{-1}^1xP_LP_N\,dx
=a_L\frac{2\delta_{N,L+1}}{2L+3}
+b_L\frac{2\delta_{N,L-1}}{2L-1}.
\]
This gives respectively
$2(L+1)/[(2L+1)(2L+3)]$ and
$2L/[(2L-1)(2L+1)]$.

\textbf{(b)} Apply the recurrence a second time:
\[
\begin{aligned}
x^2P_L
&=a_L(a_{L+1}P_{L+2}+b_{L+1}P_L)
+b_L(a_{L-1}P_L+b_{L-1}P_{L-2})\\
&=A_LP_{L+2}+B_LP_L+C_LP_{L-2},
\end{aligned}
\]
where the coefficients are
\[
A_L=\frac{(L+1)(L+2)}{(2L+1)(2L+3)},\qquad
C_L=\frac{L(L-1)}{(2L+1)(2L-1)},
\]
\[
\begin{aligned}
B_L&=\frac{(L+1)^2}{(2L+1)(2L+3)}
+\frac{L^2}{(2L+1)(2L-1)}\\
&=\frac{(L+1)^2(2L-1)+L^2(2L+3)}
{(2L+1)(2L-1)(2L+3)}\\
&=\frac{(2L+1)(2L^2+2L-1)}
{(2L+1)(2L-1)(2L+3)}
=\frac{2L^2+2L-1}{(2L-1)(2L+3)}.
\end{aligned}
\]
For $L=0,1$, direct application of $xP_0=P_1$ removes every
negative-degree term; the surviving formulas still give these
same $A_L,B_L$. Integration against $P_N$ multiplies the three
coefficients by $2/(2L+5)$, $2/(2L+1)$, and $2/(2L-3)$,
respectively. These are exactly the three stated answers.
\erratum{Both displayed case lists need a final case $0$ for every
other nonnegative integer $N$. A case with a negative value of
$N=L-1$ or $N=L-2$ is absent.}

\Needspace{5\baselineskip}
\par\bigskip
\noindent\textbf{16.3.12}\quad Since $xP_{n}(x)$ is a polynomial (of
degree $n+1$), it may be represented by the Legendre series
\[
xP_{n}(x)=\sum_{s=0}^{\infty}a_{s}P_{s}(x).
\]
\begin{enumerate}
\item[(a)] Show that $a_{s}=0$ for $s<n-1$ and $s>n+1$.
\item[(b)] \erratum{For $n\ge1$,} calculate $a_{n-1}$, $a_n$,
and $a_{n+1}$ and reproduce the recurrence relation, Eq.~(15.18).
\erratum{Handle $n=0$ separately, since $a_{-1}$ is not an
expansion coefficient.}
\end{enumerate}
\par\emph{Note.}~This argument may be put in a general form to
demonstrate the existence of a three-term recurrence relation for any
of our complete sets of orthogonal polynomials:
\[
x\varphi_{n}=a_{n+1}\varphi_{n+1}+a_{n}\varphi_{n}
+a_{n-1}\varphi_{n-1}.
\]
\solution
Since the Legendre polynomials of degrees $0,\ldots,n+1$ form a
basis of the polynomials of degree at most $n+1$, the expansion of
$xP_n$ is finite and $a_s=0$ for $s>n+1$.
For a coefficient with $s<n-1$, use the inner product:
\[
a_s=\frac{2s+1}{2}\int_{-1}^1xP_n(x)P_s(x)\,dx
=\frac{2s+1}{2}\int_{-1}^1P_n(x)[xP_s(x)]\,dx.
\]
The bracket has degree $s+1<n$. It is a combination of
$P_0,\ldots,P_{n-1}$, each orthogonal to $P_n$; hence $a_s=0$.
Only degrees $n-1,n,n+1$ can occur.

The middle coefficient vanishes because
$x[P_n(x)]^2$ is odd:
\[
a_n=\frac{2n+1}{2}\int_{-1}^1x[P_n(x)]^2dx=0.
\]
To find the upper coefficient without assuming the desired
recurrence, use Rodrigues' formula. The highest term in
$(x^2-1)^n$ is $x^{2n}$, whose $n$th derivative is
$(2n)!x^n/n!$. Thus the leading coefficient of $P_n$ is
\[
\kappa_n=\frac{(2n)!}{2^n(n!)^2}.
\]
Matching the coefficient of $x^{n+1}$ in $xP_n=a_{n+1}P_{n+1}
+a_{n-1}P_{n-1}$ gives
\[
a_{n+1}=\frac{\kappa_n}{\kappa_{n+1}}
=\frac{2(n+1)^2}{(2n+2)(2n+1)}
=\frac{n+1}{2n+1}.
\]
For $n\ge1$, evaluate at $x=1$, where every $P_j(1)=1$:
\[
1=a_{n+1}+a_{n-1},\qquad
a_{n-1}=1-\frac{n+1}{2n+1}=\frac{n}{2n+1}.
\]
Therefore
\[
\boxed{(2n+1)xP_n=(n+1)P_{n+1}+nP_{n-1}\quad(n\ge1)}.
\]
For $n=0$ the separate identity is $xP_0=P_1$.

The same degree and inner-product argument works for orthogonal
polynomials with any positive weight on their interval: multiplication
by $x$ is symmetric in that inner product, so only the neighboring
degrees and the original degree can appear. The middle coefficient
need not vanish unless the weight and interval supply the parity
symmetry used here.

\Needspace{12\baselineskip}
\subsection*{16.4 Vector Spherical Harmonics}

\Needspace{5\baselineskip}
\par\bigskip
\noindent\textbf{16.4.1}\quad Construct the vector spherical
harmonics with \erratum{total $J=0,M=0$ and total $J=1,M=0$,
for every allowed orbital degree $L$}. The notation is
$\mathbf Y_{JLM}$: \erratum{the first index is total angular momentum,
and the second is orbital angular momentum}.
\par\nobreak\hfill\emph{ANS.}\enspace
$\mathbf{Y}_{010}=-\hat{\mathbf{r}}(4\pi)^{-1/2}$\par\nobreak\hfill
$\erratum{\mathbf{Y}_{000}=\mathbf{0}\quad\text{(formal zero only)}}$\par\nobreak\hfill
$\mathbf{Y}_{120}=-\hat{\mathbf{r}}(2\pi)^{-1/2}\cos\theta
-\hat{\boldsymbol{\theta}}(8\pi)^{-1/2}\sin\theta$\par\nobreak\hfill
$\mathbf{Y}_{110}=\hat{\boldsymbol{\varphi}}\,
i\,(3/8\pi)^{1/2}\sin\theta$\par\nobreak\hfill
$\mathbf{Y}_{100}=\hat{\mathbf{r}}(4\pi)^{-1/2}\cos\theta
-\hat{\boldsymbol{\theta}}(4\pi)^{-1/2}\sin\theta$.
\solution
Use the spherical Cartesian basis
\[
\mathbf e_+=-\frac{\mathbf e_x+i\mathbf e_y}{\sqrt2},\qquad
\mathbf e_0=\mathbf e_z,\qquad
\mathbf e_-=\frac{\mathbf e_x-i\mathbf e_y}{\sqrt2}.
\]
The first index of $\mathbf Y_{JLM}$ is the total angular momentum;
the second is the orbital degree. Its definition is
\[
\mathbf Y_{JLM}=\sum_{m,q}C(L\,1\,J\mid m\,q\,M)
Y_L^m\mathbf e_q.
\]
\erratum{The triangle condition $|L-1|\le J\le L+1$ excludes
$J=L=0$. The printed $\mathbf Y_{000}=0$ is a useful formal zero,
not a normalized vector harmonic. There is also no negative-degree
orbital harmonic.}

For $J=M=0$, only $L=1$ is allowed. The singlet coefficients from
16.2.1 give
\[
\mathbf Y_{010}=\frac1{\sqrt3}
(Y_1^1\mathbf e_- -Y_1^0\mathbf e_0+Y_1^{-1}\mathbf e_+).
\]
With $\mathbf e_\rho=\cos\varphi\,\mathbf e_x+
\sin\varphi\,\mathbf e_y$, the $m=\pm1$ terms combine to
$-\sqrt{3/(4\pi)}\sin\theta\,\mathbf e_\rho$.
The $m=0$ term is
$-\sqrt{3/(4\pi)}\cos\theta\,\mathbf e_z$.
Since $\hat{\mathbf r}=\sin\theta\,\mathbf e_\rho+
\cos\theta\,\mathbf e_z$,
\[
\boxed{\mathbf Y_{010}=-\frac{\hat{\mathbf r}}{\sqrt{4\pi}}}.
\]

For $J=1,M=0,L=0$, coupling a scalar to spin one is immediate:
\[
\boxed{\mathbf Y_{100}=Y_0^0\mathbf e_0
=\frac{\cos\theta\,\hat{\mathbf r}
-\sin\theta\,\hat{\boldsymbol\theta}}{\sqrt{4\pi}}}.
\]
For $L=1$, the antisymmetric coefficients give
\[
\mathbf Y_{110}=\frac1{\sqrt2}
(Y_1^1\mathbf e_- -Y_1^{-1}\mathbf e_+).
\]
Insert $Y_1^{\pm1}=\mp\sqrt{3/(8\pi)}\sin\theta\,e^{\pm i\varphi}$.
The Cartesian bracket is
\[
-e^{i\varphi}(\mathbf e_x-i\mathbf e_y)
+e^{-i\varphi}(\mathbf e_x+i\mathbf e_y)
=2i(-\sin\varphi\,\mathbf e_x+\cos\varphi\,\mathbf e_y),
\]
so
\[
\boxed{\mathbf Y_{110}=i\sqrt{\frac3{8\pi}}\sin\theta\,
\hat{\boldsymbol\varphi}}.
\]

Finally consider $L=2$. To obtain the $J=1$ coefficients, first
write its highest $M=1$ state in the product basis
$(m,q)=(2,-1),(1,0),(0,1)$ with coefficients $(a,b,c)$.
Its raising must vanish. The coefficients of the two $M=2$
products yield
\[
\sqrt2a+2b=0,\qquad \sqrt2b+\sqrt6c=0.
\]
Thus $b=-a/\sqrt2$, $c=a/\sqrt6$, and
$a^2+b^2+c^2=5a^2/3=1$ fixes
$(a,b,c)=(\sqrt{3/5},-\sqrt{3/10},1/\sqrt{10})$.
Lowering once and dividing by $\sqrt2$ gives
\[
\begin{aligned}
C(2\,1\,1\mid1,-1,0)&=(2a+\sqrt2b)/\sqrt2=\sqrt{3/10},\\
C(2\,1\,1\mid0,0,0)&=(\sqrt6b+\sqrt2c)/\sqrt2=-\sqrt{2/5},\\
C(2\,1\,1\mid-1,1,0)&=\sqrt6c/\sqrt2=\sqrt{3/10}.
\end{aligned}
\]
Consequently
\[
\mathbf Y_{120}=\sqrt{\frac3{10}}
(Y_2^1\mathbf e_-+Y_2^{-1}\mathbf e_+)
-\sqrt{\frac25}Y_2^0\mathbf e_z.
\]
Use $Y_2^{\pm1}=\mp\sqrt{15/(8\pi)}\sin\theta\cos\theta
e^{\pm i\varphi}$ and
$Y_2^0=\sqrt{5/(16\pi)}(3\cos^2\theta-1)$:
\[
\mathbf Y_{120}=-\frac1{\sqrt{8\pi}}
\left[3\sin\theta\cos\theta\,\mathbf e_\rho+
(3\cos^2\theta-1)\mathbf e_z\right].
\]
The substitutions
$\mathbf e_\rho=\sin\theta\,\hat{\mathbf r}
+\cos\theta\,\hat{\boldsymbol\theta}$ and
$\mathbf e_z=\cos\theta\,\hat{\mathbf r}
-\sin\theta\,\hat{\boldsymbol\theta}$ give radial coefficient
$3\sin^2\theta\cos\theta+(3\cos^2\theta-1)\cos\theta=2\cos\theta$
and polar coefficient
$3\sin\theta\cos^2\theta-(3\cos^2\theta-1)\sin\theta=\sin\theta$.
Therefore
\[
\boxed{\mathbf Y_{120}
=-\frac{\cos\theta}{\sqrt{2\pi}}\hat{\mathbf r}
-\frac{\sin\theta}{\sqrt{8\pi}}\hat{\boldsymbol\theta}}.
\]
Together with the formal $\mathbf Y_{000}=0$, these reproduce
all five displayed expressions.

\Needspace{5\baselineskip}
\par\bigskip
\noindent\textbf{16.4.2}\quad \erratum{Use argument inversion
$\hat{\mathbf r}\mapsto-\hat{\mathbf r}$ with a fixed Cartesian
basis, and include only allowed index triples.} Verify that the parity of
$\mathbf{Y}_{L\,L+1\,M}$ is $(-1)^{L+1}$, that of
$\mathbf{Y}_{L\,L\,M}$ is $(-1)^{L}$, and that of
$\mathbf{Y}_{L\,L-1\,M}$ is $(-1)^{L+1}$. What happened to the
$M$-dependence of the parity?
\par\emph{Hint.}~$\hat{\mathbf{r}}$ and $\hat{\boldsymbol{\varphi}}$
have odd parity; $\hat{\boldsymbol{\theta}}$ has even parity (compare
with Exercise~3.10.25).
\solution
Under inversion of the direction,
$(\theta,\varphi)\mapsto(\pi-\theta,\varphi+\pi)$.
For integer $m$, the associated Legendre and azimuthal factors give
\[
P_l^m(-\cos\theta)=(-1)^{l+m}P_l^m(\cos\theta),\qquad
e^{im(\varphi+\pi)}=(-1)^m e^{im\varphi}.
\]
Their product is $(-1)^{l+2m}=(-1)^l$ times the original harmonic.
Thus $Y_l^m(-\hat{\mathbf r})=(-1)^lY_l^m(\hat{\mathbf r})$,
independently of $m$.

In the fixed Cartesian spherical basis, $\mathbf e_q$ is independent
of the field direction. Every term in
\[
\mathbf Y_{JLM}(\hat{\mathbf r})
=\sum_{m,q}C(L\,1\,J\mid m\,q\,M)Y_L^m(\hat{\mathbf r})\mathbf e_q
\]
has the same orbital degree $L$. Therefore
\[
\mathbf Y_{JLM}(-\hat{\mathbf r})=(-1)^L
\mathbf Y_{JLM}(\hat{\mathbf r}).
\]
Applying this to the three orbital degrees $L+1,L,L-1$ gives
\[
\boxed{\begin{aligned}
\mathbf Y_{L,L+1,M}(-\hat{\mathbf r})&=(-1)^{L+1}\mathbf Y_{L,L+1,M},\\
\mathbf Y_{L,L,M}(-\hat{\mathbf r})&=(-1)^L\mathbf Y_{L,L,M},\\
\mathbf Y_{L,L-1,M}(-\hat{\mathbf r})&=(-1)^{L+1}\mathbf Y_{L,L-1,M}.
\end{aligned}}
\]
The last case is present only when $L\ge1$.
The $M$ dependence disappears because the two scalar parity factors
contain $2m$, an even integer.

This is the argument-parity convention in the question: the
Cartesian component functions are evaluated at $-\hat{\mathbf r}$.
If one instead defines parity on a physical polar-vector field by
$(\mathcal P\mathbf F)(\mathbf r)=-\mathbf F(-\mathbf r)$,
that operation includes an additional intrinsic minus sign.
In the moving spherical basis the same calculation distributes
the argument parity among the components and the basis vectors:
$\hat{\mathbf r},\hat{\boldsymbol\varphi}$ are odd and
$\hat{\boldsymbol\theta}$ is even, as stated in the hint.

\Needspace{5\baselineskip}
\par\bigskip
\noindent\textbf{16.4.3}\quad Verify the orthonormality of the vector
spherical harmonics $\mathbf{Y}_{JLM_{J}}$
\erratum{under the Hermitian inner product
$\int\mathbf Y^*\cdot\mathbf Y'\,d\Omega$, for allowed indices}.
\solution
\erratum{The textbook's Eq.~(16.89) omits the complex conjugate on
one vector harmonic. Orthonormality uses the Hermitian product
$\int\mathbf Y^*\cdot\mathbf Y'\,d\Omega$.
An unconjugated product cannot express it: already
$\mathbf e_+\cdot\mathbf e_+=0$, whereas
$\mathbf e_+^*\cdot\mathbf e_+=1$.}
The spherical Cartesian basis obeys
$\mathbf e_q^*\cdot\mathbf e_{q'}=\delta_{qq'}$.
Write $C_{mq}^{JM}=C(L\,1\,J\mid m\,q\,M)$.
Expanding both vector harmonics before integrating gives
\[
\begin{aligned}
&\int\mathbf Y_{JLM}^*\cdot\mathbf Y_{J'L'M'}\,d\Omega\\
&\quad=\sum_{m,q}\sum_{m',q'}
(C_{mq}^{JM})^*C(L'\,1\,J'\mid m'\,q'\,M')
(\mathbf e_q^*\cdot\mathbf e_{q'})
\int Y_L^{m*}Y_{L'}^{m'}\,d\Omega\\
&\quad=\delta_{LL'}\sum_{m,q}
C(L\,1\,J\mid m\,q\,M)^*
C(L\,1\,J'\mid m\,q\,M').
\end{aligned}
\]
The remaining sum is the inner product of two normalized coupled
states expressed in the orthonormal product basis. The
Clebsch--Gordan change of basis is unitary, so the sum is
$\delta_{JJ'}\delta_{MM'}$. Therefore, for all allowed indices,
\[
\boxed{\int\mathbf Y_{JLM}^*\cdot\mathbf Y_{J'L'M'}\,d\Omega
=\delta_{JJ'}\delta_{LL'}\delta_{MM'}}.
\]
This does not assign norm one to the excluded formal zero
$\mathbf Y_{000}$.

\Needspace{5\baselineskip}
\par\bigskip
\noindent\textbf{16.4.4}\quad Jackson's \emph{Classical
Electrodynamics} (see Additional Readings) defines
$\mathbf{Y}_{LLM}$ by the equation
\[
\mathbf{Y}_{LLM}(\theta,\varphi)
=\frac{1}{\sqrt{L(L+1)}}\,\mathbf{L}\,Y_{L}^{M}(\theta,\varphi),
\]
in which the angular momentum operator $\mathbf{L}$ is given by
\[
\mathbf{L}=-i(\mathbf{r}\times\nabla).
\]
Show that this definition agrees with Eq.~(16.88)
\erratum{for $L\ge1$ and $|M|\le L$}.
\solution
\erratum{The displayed definition requires $L\ge1$. At $L=0$
its denominator vanishes, and the formal $\mathbf Y_{000}=0$
is not a normalized member of this family.}
In the spherical Cartesian basis of 16.4.1, express the vector
operator in terms of $L_\pm=L_x\pm iL_y$:
\[
\mathbf L=\mathbf e_xL_x+\mathbf e_yL_y+\mathbf e_0L_z
=\frac{\mathbf e_-L_+-\mathbf e_+L_-}{\sqrt2}+\mathbf e_0L_z.
\]
Applying the scalar ladder formulas gives
\[
\begin{aligned}
\frac{\mathbf L Y_L^M}{\sqrt{L(L+1)}}={}&
\sqrt{\frac{(L-M)(L+M+1)}{2L(L+1)}}\,Y_L^{M+1}\mathbf e_-\\
&+\frac{M}{\sqrt{L(L+1)}}Y_L^M\mathbf e_0\\
&-\sqrt{\frac{(L+M)(L-M+1)}{2L(L+1)}}\,Y_L^{M-1}\mathbf e_+.
\end{aligned}
\]
These are the three Clebsch--Gordan coefficients for coupling
orbital $L$ and vector spin $1$ to total $J=L$.

One may establish their phase and total $J$ directly. At $M=L$
the two nonzero coefficients are
$-1/\sqrt{L+1}$ on $(m,q)=(L-1,1)$ and
$\sqrt{L/(L+1)}$ on $(L,0)$. Their norm squared is one.
Raising gives the coefficient
\[
-\frac{\sqrt{2L}}{\sqrt{L+1}}
+\sqrt{\frac{L}{L+1}}\sqrt2=0
\]
on $(L,1)$. Thus this is a highest state of total $J=L$,
with positive coefficient at the largest orbital $m$.

The subsequent lowering uses the vector transformation law.
On any scalar $f$, the total lowering generator acting on
$\mathbf Lf$ satisfies
\[
J_-(\mathbf Lf)=\mathbf L(L_-f).
\]
To exhibit the cancellation, the Cartesian spin-one matrices are
$(S_i)_{jk}=-i\epsilon_{ijk}$, whereas
$[L_i,L_j]=i\epsilon_{ijk}L_k$. Therefore, for each Cartesian $i$,
\[
\begin{aligned}
J_i(\mathbf Lf)-\mathbf L(L_if)
&=\sum_{j,k}\mathbf e_j\,i\epsilon_{ijk}L_kf
+\sum_{j,k}\mathbf e_j(-i\epsilon_{ijk})L_kf=0.
\end{aligned}
\]
Taking the $x-i y$ combination proves the stated lowering identity.
It follows that lowering the displayed $M$ expression gives
$\sqrt{(L+M)(L-M+1)}$ times the expression at $M-1$.
The normalized highest state and these positive lowering factors
uniquely generate the $J=L$ multiplet. Thus its coefficients are
exactly those of Eq.~(16.88), and
\[
\boxed{\mathbf Y_{LLM}=\frac{\mathbf LY_L^M}{\sqrt{L(L+1)}}}.
\]

\Needspace{5\baselineskip}
\par\bigskip
\noindent\textbf{16.4.5}\quad \erratum{For $L\ge1$,} show that
\[
\sum_{M=-L}^{L}\mathbf{Y}_{LLM}^{*}(\theta,\varphi)\cdot
\mathbf{Y}_{LLM}(\theta,\varphi)=\frac{2L+1}{4\pi}.
\]
\par\emph{Hint.}~One way is to use Exercise~16.4.4 with $\mathbf{L}$
expanded in Cartesian coordinates and to apply raising and lowering
operators.
\solution
\erratum{This identity holds for $L\ge1$. The book's formal
$\mathbf Y_{000}=0$ would give zero on the left for $L=0$,
whereas the stated right side would be $1/(4\pi)$.}
Use 16.4.4. Since Cartesian basis vectors are orthonormal,
\[
\begin{aligned}
\sum_M\mathbf Y_{LLM}^*\cdot\mathbf Y_{LLM}
&=\frac1{L(L+1)}\sum_M
\left(|L_xY_L^M|^2+|L_yY_L^M|^2+|L_zY_L^M|^2\right)\\
&=\frac1{L(L+1)}\sum_M
\left(\frac{|L_+Y_L^M|^2+|L_-Y_L^M|^2}{2}
+M^2|Y_L^M|^2\right).
\end{aligned}
\]
The second equality follows by expanding $L_x=(L_++L_-)/2$
and $L_y=(L_+-L_-)/(2i)$: their cross terms cancel.
Now collect the coefficient of a fixed $|Y_L^k|^2$.
The raising term comes from $M=k-1$ and the lowering term
from $M=k+1$. Including the $L_z$ term, the coefficient is
\[
\begin{aligned}
&\frac12(L-k+1)(L+k)+\frac12(L+k+1)(L-k)+k^2\\
&\quad=\frac12[L(L+1)-k^2+k]
+\frac12[L(L+1)-k^2-k]+k^2\\
&\quad=L(L+1).
\end{aligned}
\]
At $k=\pm L$, the out-of-range source term has coefficient zero,
so the same collection is valid at both endpoints. Hence
\[
\sum_M\mathbf Y_{LLM}^*\cdot\mathbf Y_{LLM}
=\sum_{k=-L}^{L}|Y_L^k|^2
=\boxed{\frac{2L+1}{4\pi}},
\]
where the last step is the scalar addition theorem at coincident
directions.

\Needspace{5\baselineskip}
\par\bigskip
\noindent\textbf{16.4.6}\quad \erratum{For $L\ge1$ and $|M|\le L$,}
show that
\[
\int\erratum{\mathbf{Y}_{LLM}^{*}}\cdot
\big(\hat{\mathbf{r}}\times\mathbf{Y}_{LLM}\big)\,d\Omega=0.
\]
The integrand represents an interference term in electromagnetic
radiation that contributes to angular distributions but not to total
intensity.
\solution
The original, unconjugated version has a stronger but different
property: its integrand is identically zero.
For any vector $\mathbf A$, including one with complex components,
the bilinear scalar triple product is
\[
\mathbf A\cdot(\hat{\mathbf r}\times\mathbf A)
=\sum_{ijk}\epsilon_{ijk}A_i\hat r_jA_k=0.
\]
Interchanging $i$ and $k$ changes $\epsilon_{ijk}$ to its negative
but leaves the commuting product $A_iA_k$ unchanged. The sum
therefore equals its own negative and is zero at every point.
Taking $\mathbf A=\mathbf Y_{LLM}$ proves the original integral.

\erratum{The source has no complex conjugate in this exercise.
Accordingly its following statement about a nonzero angular
interference term cannot describe the integrand as printed.
The physically relevant Hermitian version, with the first
$\mathbf Y_{LLM}$ conjugated, can be nonzero pointwise; its
angular integral also vanishes, as shown next.}

For $L\ge1$ put $N=\sqrt{L(L+1)}$ and $Y=Y_L^M$.
Using $\mathbf L=-i\hat{\boldsymbol\varphi}\,\partial_\theta
+i\hat{\boldsymbol\theta}(\sin\theta)^{-1}\partial_\varphi$ and
$\partial_\varphi Y=iMY$ gives
\[
A_\theta=-\frac{M}{N\sin\theta}Y,\qquad
A_\varphi=-\frac{i}{N}\partial_\theta Y,\qquad A_r=0.
\]
Since $\hat{\mathbf r}\times\hat{\boldsymbol\theta}
=\hat{\boldsymbol\varphi}$ and
$\hat{\mathbf r}\times\hat{\boldsymbol\varphi}
=-\hat{\boldsymbol\theta}$,
\[
\begin{aligned}
\mathbf A^*\cdot(\hat{\mathbf r}\times\mathbf A)
&=-A_\theta^*A_\varphi+A_\varphi^*A_\theta\\
&=-\frac{iM}{L(L+1)\sin\theta}
\left[Y^*\partial_\theta Y+
(\partial_\theta Y^*)Y\right]\\
&=-\frac{iM}{L(L+1)\sin\theta}\partial_\theta|Y|^2.
\end{aligned}
\]
The modulus is independent of $\varphi$, so integration gives
\[
\int\mathbf A^*\cdot(\hat{\mathbf r}\times\mathbf A)\,d\Omega
=-\frac{2\pi iM}{L(L+1)}
\left[|Y_L^M(\theta,0)|^2\right]_{\theta=0}^{\theta=\pi}=0.
\]
For $M=0$ the prefactor is zero; for $M\ne0$ the spherical
harmonic vanishes at both poles. For example $L=M=1$ gives
$|Y_1^1|^2=3\sin^2\theta/(8\pi)$ and hence
\[
-\frac{i}{2\sin\theta}\frac{d}{d\theta}
\left(\frac{3\sin^2\theta}{8\pi}\right)
=-\frac{3i}{8\pi}\cos\theta,
\]
which is nonzero away from the equator but integrates to zero. Thus the corrected
Hermitian interpretation has the angular behavior described
in the prose.

